% =============================================================
% YSU Master's Thesis — Gurgen Blbulyan style
% Compiler: XeLaTeX
% Author: Nerses Asaturyan
% Topic: Electronic Voting Systems
% =============================================================
\documentclass[12pt,a4paper,oneside]{report}

% ---------- ENCODING & FONTS ----------
\usepackage{fontspec}
\usepackage{polyglossia}
\setdefaultlanguage{english}
\setotherlanguage{armenian}

% Main font — must support Armenian Unicode.
% On Overleaf, "FreeSerif" works out of the box. 
% For nicer look, upload "GHEA Grapalat" or "Sylfaen" and replace below.
\setmainfont{FreeSerif}
\newfontfamily\armenianfont[Script=Armenian]{FreeSerif}

% ---------- PAGE GEOMETRY (per YSU regulation §20) ----------
\usepackage[a4paper,
            left=3.0cm, right=1.5cm,
            top=2.0cm,  bottom=2.5cm]{geometry}

% ---------- LINE SPACING ----------
\usepackage{setspace}
\onehalfspacing  % 1.5 line spacing

% ---------- PARAGRAPH ----------
\setlength{\parindent}{1cm}   % 1 cm first-line indent (§20)
\setlength{\parskip}{0pt}

% ---------- PACKAGES ----------
\usepackage{graphicx}
% Build runs from thesis/build/, so paths are relative to that dir.
\graphicspath{{../results/plots/}{../results/plots/criterion/}}

% Per-figure data-source line, kept short and consistent.
\newcommand{\figdata}[1]{\\\textit{Data:} \texttt{#1}.}
\usepackage{amsmath, amssymb, amsthm}
\usepackage{mathtools}
\usepackage{float}
\usepackage{caption}
\usepackage{subcaption}
\usepackage{booktabs}
\usepackage{array}
\usepackage{longtable}
\usepackage{enumitem}
\usepackage{hyperref}
\usepackage{url}
\usepackage{xcolor}
\usepackage{fancyhdr}
\usepackage{titlesec}
\usepackage{tocloft}
\usepackage{indentfirst}

% ---------- HYPERREF ----------
\hypersetup{
  colorlinks=true,
  linkcolor=black,
  urlcolor=blue,
  citecolor=black,
  pdftitle={Electronic Voting Systems},
  pdfauthor={Nerses Asaturyan}
}

% ---------- CAPTIONS (10pt, centered, per §21) ----------
\captionsetup{font=small, labelfont=it, textfont=it, justification=centering}

% ---------- CHAPTER / SECTION HEADINGS ----------
% Chapter titles: bold, uppercase, new page (§20)
\titleformat{\chapter}[block]
  {\normalfont\fontsize{14}{16}\bfseries\filcenter}
  {\thechapter.}{1em}{\MakeUppercase}
\titlespacing*{\chapter}{0pt}{0pt}{20pt}

\titleformat{\section}
  {\normalfont\fontsize{12}{14}\bfseries}
  {\thesection}{1em}{}
\titleformat{\subsection}
  {\normalfont\fontsize{12}{14}\bfseries\itshape}
  {\thesubsection}{1em}{}

% ---------- PAGE NUMBERING (centered at bottom, per §20) ----------
\pagestyle{plain}

% ---------- BIBLIOGRAPHY ----------
\usepackage[numbers,square,sort&compress]{natbib}
\bibliographystyle{unsrt}

% =============================================================
\begin{document}

% =============================================================
% TITLE PAGE — HAVELVAĐ 1 (Armenian)
% =============================================================
\begin{titlepage}
\centering
\begin{armenian}
\fontsize{18}{22}\selectfont\textbf{ԵՐԵՎԱՆԻ ՊԵՏԱԿԱՆ ՀԱՄԱԼՍԱՐԱՆ}
\end{armenian}
\\[1.2cm]
\begin{armenian}
\fontsize{18}{22}\selectfont\textbf{ՄԱԹԵՄԱՏԻԿԱՅԻ ԵՎ ՄԵԽԱՆԻԿԱՅԻ ՖԱԿՈՒԼՏԵՏ}
\end{armenian}
\\[0.8cm]
\begin{armenian}
\fontsize{16}{20}\selectfont\textbf{ԱԿՏՈՒԱՐԱԿԱՆ ԵՎ ՖԻՆԱՆՍԱԿԱՆ ՄԱԹԵՄԱՏԻԿԱՅԻ ԱՄԲԻՈՆ}
\end{armenian}
\\[0.8cm]
\begin{armenian}
\fontsize{18}{22}\selectfont\textbf{ԲԼՈԿՉԵՅՆ ԵՎ ԹՎԱՅԻՆ ԱՐԺՈՒՅԹՆԵՐ ՄԱԳԻՍՏՐՈՍԱԿԱՆ ԿՐԹԱԿԱՆ ԾՐԱԳԻՐ}
\end{armenian}
\\[2.5cm]
\begin{armenian}
\fontsize{18}{22}\selectfont\textbf{ԱՍԱՏՈՒՐՅԱՆ ՆԵՐՍԵՍ ԱՐՏԱՎԱԶԴԻ}
\end{armenian}
\\[1.5cm]
\begin{armenian}
\fontsize{18}{22}\selectfont\textbf{ՄԱԳԻՍՏՐՈՍԱԿԱՆ ԹԵԶ}
\end{armenian}
\\[0.8cm]
\begin{armenian}
\fontsize{16}{20}\selectfont\textbf{ԷԼԵԿՏՐՈՆԱՅԻՆ ՔՎԵԱՐԿՈՒԹՅԱՆ ՀԱՄԱԿԱՐԳԵՐ}
\end{armenian}
\\[1.5cm]
\begin{armenian}
\fontsize{14}{18}\selectfont\textit{\textbf{«Բլոկչեյն և թվային արժույթներ» մասնագիտացմամբ\\
ֆինանսական մաթեմատիկայի մագիստրոսի որակավորման\\
աստիճանի հայցման համար}}
\end{armenian}

\vfill

\begin{armenian}
\fontsize{13}{16}\selectfont\textbf{ԵՐԵՎԱՆ 2026}
\end{armenian}
\end{titlepage}

% =============================================================
% SIGNATURES PAGE — HAVELVAĐ 2 (Armenian)
% =============================================================
\thispagestyle{empty}
\vspace*{2cm}

\noindent
\begin{minipage}[t]{0.28\textwidth}
\begin{armenian}\fontsize{13}{16}\selectfont\textit{\textbf{Ուսանող`}}\end{armenian}
\end{minipage}%
\begin{minipage}[t]{0.6\textwidth}
\vspace*{0.7cm}\par
\centering\begin{armenian}\small Ասատուրյան Ներսես Արտավազդի\end{armenian}
\end{minipage}\\[1.5cm]

\noindent
\begin{minipage}[t]{0.28\textwidth}
\begin{armenian}\fontsize{13}{16}\selectfont\textit{\textbf{Գիտական ղեկավար`}}\end{armenian}
\end{minipage}%
\begin{minipage}[t]{0.6\textwidth}
\vspace*{0.7cm}\par
\centering\begin{armenian}\small ֆ.մ.գ.թ. Դավիդովա Դիանա\end{armenian}
\end{minipage}\\[1.5cm]

\noindent\begin{armenian}\fontsize{13}{16}\selectfont\textbf{«Թույլատրել պաշտպանության»}\end{armenian}\\[1cm]

\noindent
\begin{minipage}[t]{0.28\textwidth}
\begin{armenian}\fontsize{13}{16}\selectfont\textit{\textbf{Ամբիոնի վարիչ`}}\end{armenian}
\end{minipage}%
\begin{minipage}[t]{0.6\textwidth}
\vspace*{0.7cm}\par
\centering\begin{armenian}\small տնտեսագիտության թեկնածու, դոցենտ Ավետիսյան Կարինե\end{armenian}
\end{minipage}\\[2cm]

\noindent\begin{armenian}«\rule{1cm}{0.4pt}»\rule{3cm}{0.4pt}20\rule{1cm}{0.4pt}թ.\end{armenian}
\newpage

% =============================================================
% ABSTRACT / SUMMARY (trilingual title block: Armenian, English, Russian)
% =============================================================
\thispagestyle{empty}
\begin{center}
\begin{armenian}
{\large\textbf{Էլեկտրոնային քվեարկության համակարգեր}}\\[0.4cm]
\end{armenian}
{\large\textbf{Electronic Voting Systems}}\\[0.4cm]
{\large\textbf{Системы электронного голосования}}\\[1.2cm]

\textbf{ABSTRACT}
\end{center}
\vspace{0.4cm}

\noindent\textbf{Goal.} The goal of this work is to study modern electronic 
voting systems, analyze in detail the Aura private voting protocol 
(Jivanyan, Feickert, 2022), and deliver its first open-source 
implementation with practical performance benchmarks. The work aims 
to demonstrate that the Aura protocol is practically applicable to 
digital direct-democracy systems.\\[0.3cm]

\noindent\textbf{Results.} The first open-source implementation of the 
Aura protocol has been developed and is publicly available on GitHub. 
Every phase of the implementation --- distributed key generation, voter 
registration, encryption, zero-knowledge proofs, and threshold decryption 
--- has been benchmarked under varying parameters. A step-by-step 
numerical example of the Aura protocol is also included to make the 
abstract construction concrete.

\vspace{0.5cm}
\noindent\textbf{Code is available at:}\\
\url{https://github.com/nerses-asaturyan/aura}

\newpage

% =============================================================
% TABLE OF CONTENTS (counts as page 4 per §20)
% =============================================================
\setcounter{page}{4}
\tableofcontents
\newpage

% =============================================================
% INTRODUCTION
% =============================================================
\chapter*{Introduction}
\addcontentsline{toc}{chapter}{Introduction}

\section*{Motivation}
Democratic elections are one of the most fundamental institutions of modern 
society, yet the process of casting and counting votes still relies heavily 
on paper ballots, trusted intermediaries, and centralized election authorities. 
As digital infrastructure becomes ubiquitous, there is growing interest in 
moving elections --- and more generally, collective decision-making processes 
--- onto cryptographic protocols that can be verified by anyone, while still 
protecting the privacy of individual voters.

Electronic voting (e-voting) systems aim to provide three guarantees that 
are notoriously hard to achieve simultaneously: \textbf{ballot privacy} 
(no one learns how a specific voter voted), \textbf{verifiability} (anyone 
can check that the tally is correct), and \textbf{eligibility} (only 
registered voters can vote, and each voter can vote at most once). Modern 
cryptographic e-voting schemes attempt to satisfy all three using building 
blocks like homomorphic encryption, zero-knowledge proofs, threshold 
cryptography, and blockchain-based bulletin boards.

Beyond traditional government elections, the same machinery enables a much 
broader vision: \textbf{direct digital democracy}, in which citizens vote 
frequently on individual policy questions rather than electing representatives 
who decide on their behalf. For such systems to be trusted at scale, they 
must be both cryptographically sound and practically efficient.

\section*{Goal and Objectives}
The goal of this thesis is to study modern cryptographic e-voting systems, 
with a focus on the \textit{Aura private voting protocol} of Jivanyan and 
Feickert (2022)~\cite{aura2022}, and to deliver the first open-source 
implementation of the protocol along with practical performance benchmarks.

The specific objectives are:
\begin{itemize}
    \item Survey existing cryptographic e-voting systems briefly to position 
    Aura among them.
    \item Provide a self-contained presentation of the cryptographic primitives 
    used in Aura (ElGamal encryption, Pedersen commitments, Merkle trees, 
    Schnorr proofs, one-of-many proofs, threshold secret sharing).
    \item Describe the full Aura protocol from setup through tally, with 
    a step-by-step worked numerical example.
    \item Implement the protocol end-to-end in code and publish it as 
    an open-source repository.
    \item Benchmark each component of the protocol under varying parameters 
    (number of voters, candidates, threshold size) to evaluate practical viability.
\end{itemize}

\section*{Object and Subject of Study}
The \textbf{object} of study is cryptographic electronic voting systems. 
The \textbf{subject} is the Aura protocol --- its construction, security 
guarantees, and concrete computational cost when implemented.

\section*{Hypothesis}
The Aura protocol, despite combining several heavyweight cryptographic 
primitives (threshold ElGamal, one-of-many proofs, range proofs), is 
practically efficient enough to support real-world voting scenarios --- 
including frequent direct-democracy referenda --- on commodity hardware.

\section*{Theoretical and Methodological Foundations}
The work draws on results from public-key cryptography, zero-knowledge 
proof systems, and distributed computing. Methodologically, the thesis 
combines (i) a brief review of existing e-voting schemes, (ii) a formal 
exposition of the Aura protocol, (iii) a software-engineering effort to 
implement the protocol, and (iv) empirical performance measurement.

\section*{Scientific and Practical Significance}
\textbf{Scientifically}, this thesis contributes the first open and 
reproducible implementation of the Aura protocol, enabling other researchers 
to verify the authors' claims, build on the codebase, or compare against 
new schemes. It also includes a worked numerical example of Aura that 
makes the protocol's mechanics easier to follow.

\textbf{Practically}, the benchmarks provide concrete guidance to anyone 
considering deploying Aura --- for example, in DAO governance, university 
student elections, or municipal direct-democracy pilots --- by quantifying 
the cost of each protocol step.

\section*{Thesis Structure}
The remainder of this thesis is organized as follows. 
Chapter~\ref{ch:literature} briefly reviews related work in cryptographic 
e-voting. Chapter~\ref{ch:primitives} introduces the cryptographic 
primitives required to understand Aura. Chapter~\ref{ch:aura} describes 
the Aura protocol in full. Chapter~\ref{ch:byhand} walks through the 
protocol with a numerical example. Chapter~\ref{ch:implementation} 
documents the open-source implementation. Chapter~\ref{ch:results} 
reports benchmark results. Chapter~\ref{ch:discussion} discusses 
limitations and future work, and the final chapter concludes.

% =============================================================
% CHAPTER 1 — RELATED WORK (brief)
% =============================================================
\chapter{Related Work}
\label{ch:literature}

Cryptographic e-voting has been studied for several decades, and several 
practical systems have been deployed at varying scales. Helios~\cite{helios} 
is among the earliest deployed cryptographic e-voting systems and uses 
ElGamal encryption with cast-or-audit ballot verification; it has been 
used for university and IACR elections, although it relies on a small 
set of trusted election servers and does not provide coercion resistance. 
The Russian federal remote e-voting scheme deployed in 2021~\cite{russian2021} 
is a more recent large-scale deployment built on a blockchain backend and 
blind signatures, but has been criticized for limited transparency and 
verifiability. Broader surveys of the field 
include~\cite{cryptoreview} on the cryptographic side and~\cite{blockchainreview} 
on the blockchain-based side, both of which categorize existing systems 
into mix-net-based, homomorphic-tally-based, and blind-signature-based 
families.

Aura~\cite{aura2022} was introduced in 2022 by Jivanyan and Feickert as a 
private voting protocol with reduced trust assumptions on the tallying 
authorities. Compared to Helios, it adds threshold decryption and 
Groth-Kohlweiss one-of-many proofs for stronger anonymity; compared to 
most blockchain-based schemes, it does not require a smart-contract 
platform --- only a public bulletin board. To date, no public 
implementation of Aura exists, which motivates the engineering 
contribution of this thesis.

% =============================================================
% CHAPTER 2 — CRYPTOGRAPHIC PRIMITIVES
% =============================================================
\chapter{Cryptographic Primitives}
\label{ch:primitives}

Aura is built from several standard cryptographic primitives. This chapter 
gives a self-contained introduction to each, sufficient to follow the 
protocol description in Chapter~\ref{ch:aura}. Readers familiar with 
public-key cryptography may skim this chapter.

\section{Cyclic Groups and the Discrete Logarithm Problem}
\label{sec:cyclic-groups}

The cryptographic primitives used in Aura are all built on top of a finite 
cyclic group in which the discrete logarithm problem is believed to be 
intractable. Throughout this thesis, $\mathbb{G}$ denotes a cyclic group 
of prime order $q$ with generator $g$, written multiplicatively, so that
\[
    \mathbb{G} = \{ g^0, g^1, g^2, \ldots, g^{q-1} \},
\]
and exponents live in the finite field $\mathbb{Z}_q$. Working in a 
prime-order group means every non-identity element is itself a generator 
and admits cleaner security proofs.

Given $g$ and $h = g^x$, the \textit{discrete logarithm problem} (DLP) 
is to recover $x \in \mathbb{Z}_q$. Computing $h$ from $x$ takes only 
$O(\log q)$ group multiplications, but the reverse direction is believed 
to be hard when $q$ is large; the best known generic algorithms run in 
time roughly $\sqrt{q}$, which is infeasible once $q$ has on the order of 
$256$ bits. This is what makes $g^x$ usable as a one-way commitment to 
$x$ --- a property Aura relies on, most directly in the serial number 
$\text{sn}_i = u^{s_i}$.

A stronger but related hardness assumption, used by ElGamal encryption 
in Aura, is the \textit{Decisional Diffie-Hellman} (DDH) assumption: 
given $g, g^a, g^b$ for random $a, b \in \mathbb{Z}_q$, no efficient 
algorithm can distinguish $g^{ab}$ from a uniformly random element of 
$\mathbb{G}$.

In practice, $\mathbb{G}$ is realized either as a prime-order subgroup 
of $\mathbb{Z}_p^*$ (the setting of the worked example in 
Chapter~\ref{ch:byhand}, with toy parameters $p = 47$, $q = 23$) or 
as an elliptic-curve group (the setting of the implementation in 
Chapter~\ref{ch:implementation}). Elliptic-curve groups achieve the 
same security level at much smaller key sizes, which is why every 
modern deployment of discrete-log-based cryptography uses them.
\section{ElGamal Encryption}
\label{sec:elgamal}

ElGamal encryption~\cite{elgamal} is a public-key cryptosystem whose security 
reduces to the DDH assumption in $\mathbb{G}$. Aura uses it as the encryption 
layer for every ballot.

\subsection*{Standard ElGamal}
The secret key is a uniformly random $sk \in \mathbb{Z}_q$, and the public 
key is $pk = g^{sk} \in \mathbb{G}$. To encrypt a message $m \in \mathbb{G}$ 
with fresh randomness $r \in \mathbb{Z}_q$, the sender computes
\[
    \text{Enc}_{pk}(m; r) = (c_1, c_2) = (g^r,\; m \cdot pk^r).
\]
Decryption recovers $m$ as $m = c_2 / c_1^{sk}$, since 
$c_1^{sk} = g^{r \cdot sk} = pk^r$. ElGamal is \textit{multiplicatively 
homomorphic}: the componentwise product of two ciphertexts encrypts the 
product of the corresponding plaintexts,
\[
    \text{Enc}_{pk}(m_1) \cdot \text{Enc}_{pk}(m_2) = \text{Enc}_{pk}(m_1 \cdot m_2).
\]

\subsection*{Exponential ElGamal}
For voting, an additive homomorphism is needed, so that ciphertexts of $0/1$ 
votes can be multiplied to encrypt their sum. \textit{Exponential ElGamal} 
achieves this by encoding an integer message $v \in \mathbb{Z}_q$ as the 
group element $g^v$ before encryption:
\[
    \text{Enc}_{pk}(v; r) = (g^r,\; g^v \cdot pk^r).
\]
The componentwise product of two such ciphertexts now encrypts $g^{v_1 + v_2}$, 
turning the multiplicative homomorphism into an additive one. The cost is 
that decryption recovers only $g^v$, not $v$ itself --- but when $v$ is 
known to lie in a small range (such as a vote tally $T_j \in [0, n]$), $v$ 
can be read off by lookup against a precomputed table $\{g^0, g^1, \ldots, g^n\}$.

This is the variant Aura uses: each ballot slot encrypts a bit $v_{i,j} \in \{0, 1\}$ 
as $E_{i,j} = (g^{\rho_{i,j}}, g^{v_{i,j}} \cdot pk^{\rho_{i,j}})$, and the 
homomorphic column product $E_\Sigma^{(j)} = \prod_i E_{i,j}$ then encrypts 
the candidate total $T_j = \sum_i v_{i,j}$.
\section{Pedersen Commitments}
\label{sec:pedersen}

A \textit{commitment scheme} lets a party publish a short value that 
binds them to a secret without revealing it, and later open the commitment 
to prove what was committed. Pedersen commitments~\cite{pedersen} are the 
flavor used in Aura, because they are both \textit{perfectly hiding} and 
\textit{computationally binding} under the discrete logarithm assumption.

\subsection*{Definition}
Fix two generators $g, h \in \mathbb{G}$ such that no party knows the 
discrete logarithm of $h$ with respect to $g$. To commit to a value 
$v \in \mathbb{Z}_q$, the committer samples a fresh random 
$r \in \mathbb{Z}_q$ (the \textit{blinding factor}) and publishes
\[
    \text{cm}(v, r) = g^v h^r \in \mathbb{G}.
\]
To open the commitment later, the committer reveals $(v, r)$, and any 
verifier checks that $g^v h^r$ matches the published $\text{cm}(v, r)$.

\subsection*{Security Properties}
\textbf{Perfectly hiding.} Because $r$ is uniform in $\mathbb{Z}_q$ and 
multiplication by $h^r$ is a bijection on $\mathbb{G}$, the commitment 
$\text{cm}(v, r)$ is uniformly distributed in $\mathbb{G}$ regardless of $v$. 
No adversary, even with unbounded computing power, learns anything about 
$v$ from $\text{cm}(v, r)$.

\textbf{Computationally binding.} Suppose a committer could open the same 
commitment to two distinct pairs $(v, r) \ne (v', r')$. Then 
$g^v h^r = g^{v'} h^{r'}$, which rearranges to $h = g^{(v - v') / (r' - r)}$, 
revealing the discrete logarithm of $h$ with respect to $g$. Under the DLP 
assumption (Section~\ref{sec:cyclic-groups}), this is computationally 
infeasible, so the committer is bound to her original choice of $v$.

These two properties trade off against each other: hiding is unconditional 
(perfect), while binding holds only against computationally bounded 
adversaries. This trade-off is fundamental --- no commitment scheme can 
be both perfectly hiding and perfectly binding at once.

\subsection*{Why Aura Uses Pedersen Commitments}
In Phase 2 of the protocol, each voter $V_i$ registers by publishing 
$\text{cm}_i = g^{s_i} h^{r_i}$, where $s_i$ is her private seed and 
$r_i$ is fresh randomness. Both properties of Pedersen commitments are 
essential:

\begin{itemize}
    \item \textit{Perfect hiding} guarantees that the public registration 
    list $L = (\text{cm}_1, \ldots, \text{cm}_n)$ leaks no information 
    about any voter's seed --- even an attacker with unbounded resources 
    cannot link a commitment to a specific $s_i$.
    \item \textit{Computational binding} guarantees that a voter cannot 
    later claim a different seed than the one she registered with, which 
    would let her vote multiple times under different identities or 
    forge eligibility proofs.
\end{itemize}

The commitment is also \textit{additively homomorphic}: 
$\text{cm}(v_1, r_1) \cdot \text{cm}(v_2, r_2) = \text{cm}(v_1 + v_2, r_1 + r_2)$. 
Aura does not use this property in the registration phase, but it is 
exploited inside the Groth-Kohlweiss one-of-many proof 
(Section~\ref{sec:one-of-many}) that proves a voter's commitment lies 
in $L$ without revealing which one.
\section{Merkle Trees}
\label{sec:merkle-trees}

A \textit{Merkle tree}~\cite{merkle} is a hash tree that compresses a list 
of values into a single short hash --- the \textit{root} --- in such a way 
that membership of any individual value can be proven with only 
$O(\log n)$ data. Aura itself does not use a Merkle tree, but the 
construction is included here as standard background, and because the 
worked numerical example in Chapter~\ref{ch:byhand} introduces one as a 
pedagogical simplification.

\subsection*{Construction}
Let $H_{\text{MT}}$ be a collision-resistant hash function. Given a list 
of $n$ leaves $(\ell_1, \ldots, \ell_n)$, the tree is built bottom-up: 
each internal node is the hash of its children,
\[
    \text{parent} = H_{\text{MT}}(\text{children}),
\]
and the topmost node is the root, denoted $\text{MT.root}$. The root is 
a short ($O(1)$ size) commitment to the full list: changing any leaf, 
anywhere in the tree, changes the root with overwhelming probability 
under collision resistance.

\subsection*{Authentication Paths}
To prove that a particular leaf $\ell_i$ is in the tree without revealing 
the other leaves, the prover supplies the \textit{authentication path}: 
the sibling hashes encountered on the way from $\ell_i$ up to the root. 
A verifier who knows $\text{MT.root}$ recomputes the path by hashing 
$\ell_i$ together with each set of siblings in turn, and checks that the 
computed root matches the published one. The path has length 
$O(\log n)$, so membership proofs take $O(\log n)$ space and time.

\subsection*{Why Aura Does Not Use One}
A natural intuition is that Merkle trees are necessary to make voter-list 
membership proofs scale: with $n$ registered voters, a flat list takes 
$O(n)$ space, while a Merkle root takes $O(1)$. However, the 
Groth-Kohlweiss commitment-set proof (Section~\ref{sec:one-of-many}) 
already produces $O(\log n)$-size membership proofs operating directly 
on a flat list of commitments --- without any tree structure on top of it. 
Aura therefore publishes the full voter commitment list 
$L = (\text{cm}_1, \ldots, \text{cm}_n)$ and runs the commitment-set 
proof on $L$ directly, achieving logarithmic proof size without the 
extra Merkle layer. The flat-list approach also avoids the bookkeeping 
of authentication paths that voters would otherwise need to keep 
synchronized with a published root.
\section{Schnorr Identification and Sigma Protocols}
\label{sec:schnorr}

A \textit{sigma protocol}~\cite{cramer-damgard} is a three-move interactive 
proof in which a prover convinces a verifier that she knows a witness $w$ 
for a public statement $y$ --- without revealing anything about $w$ --- by 
exchanging a \textit{commitment}, a \textit{challenge}, and a \textit{response}. 
Schnorr's protocol~\cite{schnorr} is the canonical example: it proves 
knowledge of a discrete logarithm.

\subsection*{Schnorr's Protocol}
The prover $P$ knows $x \in \mathbb{Z}_q$, the verifier $V$ knows the 
statement $y = g^x \in \mathbb{G}$, and they wish to convince $V$ that 
$P$ knows $x$. The interaction is:
\begin{enumerate}
    \item $P$ samples $r \leftarrow \mathbb{Z}_q$, sends the commitment 
    $t = g^r$.
    \item $V$ samples a random challenge $c \leftarrow \mathbb{Z}_q$ and 
    sends it.
    \item $P$ replies with the response $z = r + c x \mod q$.
\end{enumerate}
The verifier accepts iff $g^z = t \cdot y^c$, which holds for an honest 
prover since $g^z = g^{r + cx} = t \cdot y^c$.

\subsection*{Security Properties}
Schnorr's protocol satisfies two properties that together define a sigma 
protocol. \textit{Special soundness:} from any two accepting transcripts 
$(t, c, z)$ and $(t, c', z')$ with the same commitment but different 
challenges, the witness can be extracted as $x = (z - z') / (c - c') \mod q$. 
A cheating prover that succeeds non-negligibly could be rewound to recover 
$x$, contradicting the DLP. \textit{Honest-verifier zero knowledge:} a 
simulator can produce transcripts indistinguishable from real ones by 
picking $c, z$ uniformly and setting $t = g^z \cdot y^{-c}$, so the 
protocol leaks nothing about $x$.

\subsection*{The Fiat-Shamir Transform}
The protocol above is interactive. The \textit{Fiat-Shamir transform}~\cite{fiatshamir} 
makes it non-interactive by replacing the verifier's challenge with a 
hash of the statement and the commitment,
\[
    c = H(y,\; t),
\]
where $H$ is modeled as a random oracle (Section~\ref{sec:random-oracle}). 
The prover produces the entire transcript on her own and sends $(t, z)$; 
any verifier with $y$ recomputes $c = H(y, t)$ and checks 
$g^z = t \cdot y^c$. Soundness in the random oracle model rests on a 
specific property: because the prover must commit to $t$ before seeing 
$c$, and because $H$ behaves as a truly random function, she cannot 
manipulate $t$ to make a particular $c$ come out of the hash --- so the 
hash output is, from her perspective, indistinguishable from a 
freshly-drawn random challenge. This is the construction used throughout 
Aura: every interactive proof in the protocol is a sigma protocol made 
non-interactive in exactly this way.
\section{Disjunctive Proofs (OR Proofs)}
\label{sec:or-proofs}

Sigma protocols prove a single statement at a time. For voting, what is 
needed is a proof that \textit{one of two} statements is true, without 
revealing which one. The \textit{Cramer-Damg\aa rd-Schoenmakers 
transformation}~\cite{cramer-damgard} composes any two sigma protocols 
into such a disjunctive (OR) proof.

\subsection*{Construction}
Suppose the prover wants to prove $\text{Stmt}_0 \lor \text{Stmt}_1$ 
knowing a witness only for the true branch $b \in \{0, 1\}$. The trick is 
to use the simulator (from honest-verifier zero knowledge in 
Section~\ref{sec:schnorr}) to fake the false branch, while running the 
real protocol on the true one:
\begin{enumerate}
    \item For the false branch $1 - b$, pick a random $c_{1-b}$ and run 
    the simulator to produce $(t_{1-b}, c_{1-b}, z_{1-b})$.
    \item For the true branch $b$, compute the real commitment $t_b$.
    \item Send $(t_0, t_1)$; the verifier replies with a single challenge 
    $c$.
    \item Set $c_b = c - c_{1-b}$ and compute the real response $z_b$.
    \item Send $(c_0, c_1, z_0, z_1)$.
\end{enumerate}
The verifier checks that $c_0 + c_1 = c$ and that both sub-transcripts 
verify. The split of $c$ into $(c_0, c_1)$ binds the two branches: the 
prover can pre-pick one sub-challenge, but the other is forced by $c$, 
so she can only complete the protocol on a branch for which she has a 
witness. Both branches verify regardless of $b$, so the verifier learns 
nothing about which one was real. After Fiat-Shamir, with 
$c = H(y_0, y_1, t_0, t_1)$, the whole proof is a single message.

\subsection*{Use in Aura: Bit Proofs}
Each ballot ciphertext slot $E_{i,j} = (c_{1,j}, c_{2,j})$ should encrypt 
either $0$ or $1$, but a malicious voter could try to encrypt a larger 
value to give a candidate extra votes. The \textit{bit proof} forces 
$v_{i,j} \in \{0, 1\}$ via an OR of two Schnorr-style statements:
\[
    \underbrace{c_{2,j} = pk^{\rho_{i,j}}}_{v_{i,j} = 0}
    \quad \text{OR} \quad
    \underbrace{c_{2,j} / g = pk^{\rho_{i,j}}}_{v_{i,j} = 1}.
\]
The voter knows $\rho_{i,j}$ for exactly one branch (the one matching 
her actual bit) and uses the OR construction above to produce a proof 
that verifies for both --- without revealing which.
\section{One-of-Many Proofs (Groth-Kohlweiss)}
\label{sec:one-of-many}

The OR proof of Section~\ref{sec:or-proofs} extends to disjunctions of $n$ 
statements, but at a cost: proof size grows linearly in $n$. For an Aura 
election with thousands of voters, this is infeasible. The 
\textit{Groth-Kohlweiss one-of-many proof}~\cite{groth-kohlweiss} solves 
this: it proves that one of $n$ public commitments opens to zero, in 
proof size $O(\log n)$, without revealing which.

\subsection*{Construction Sketch}
The prover holds commitments $(\text{cm}_1, \ldots, \text{cm}_n)$ and 
an index $\ell$ such that she knows an opening of $\text{cm}_\ell$ to 
zero. Instead of running $n$ disjunctive branches, she commits to the 
bits of $\ell$ in binary, $\ell = (\ell_1, \ldots, \ell_m)$ with 
$m = \lceil \log_2 n \rceil$, and runs $m$ parallel Schnorr-style proofs 
that each bit commitment opens to $0$ or $1$. The challenge-response 
phase mixes these bit commitments with the full list in a way that 
algebraically cancels every term except the $\ell$-th, leaving exactly 
the statement that $\text{cm}_\ell$ opens to zero. The prover transmits 
only $O(\log n)$ group elements and scalars in total.

\subsection*{Use in Aura}
At ballot construction, each voter must prove she is one of the registered 
voters in $L = (\text{cm}_1, \ldots, \text{cm}_n)$ without revealing her 
index. Aura applies the Groth-Kohlweiss proof to a shifted list of 
commitments encoding her seed $s_i$ and serial number 
$\text{sn}_i = u^{s_i}$, so that the same proof simultaneously establishes 
eligibility and consistency between $\text{cm}_i$ and $\text{sn}_i$. This 
is the core anonymity primitive of the protocol: it is what keeps ballots 
small while still hiding the voter's identity.
\section{Shamir Secret Sharing and Threshold Cryptography}
\label{sec:threshold}

Shamir's secret sharing~\cite{shamir} splits a secret into $n$ shares 
such that any $t$ shares can reconstruct the secret, but fewer than $t$ 
shares reveal nothing about it. \textit{Threshold cryptography} extends 
this idea so that the shareholders can jointly use the secret --- for 
instance, to decrypt a ciphertext --- without ever assembling it in one 
place. Aura uses both techniques to distribute the ElGamal decryption 
key across $\eta$ tallying authorities.

\subsection*{Shamir's Scheme}
To share a secret $s \in \mathbb{Z}_q$ with threshold $t$, the dealer 
samples a random polynomial of degree $t - 1$ over $\mathbb{Z}_q$ with 
constant term $s$:
\[
    f(x) = s + a_1 x + a_2 x^2 + \cdots + a_{t-1} x^{t-1}.
\]
Share $i$ is $f(i)$, given to party $i$. Any $t$ shares determine $f$ 
uniquely by Lagrange interpolation, recovering $s = f(0)$. With fewer 
than $t$ shares, $s$ is information-theoretically hidden: every value 
of $s$ is consistent with the missing shares.

\subsection*{Verifiable Secret Sharing}
Plain Shamir sharing assumes an honest dealer. \textit{Feldman verifiable 
secret sharing}~\cite{feldman} adds public commitments to the polynomial 
coefficients: the dealer publishes $C_k = g^{a_k}$ for $k = 0, \ldots, t - 1$ 
(with $a_0 = s$), allowing each shareholder $i$ to verify that her share 
satisfies
\[
    g^{f(i)} \stackrel{?}{=} \prod_{k = 0}^{t-1} (C_k)^{i^k}.
\]
A cheating dealer cannot distribute inconsistent shares without being 
caught. The trade-off is that $C_0 = g^s$ is now public, which leaks 
$g^s$ but not $s$ itself --- exactly what is needed when $g^s$ will be 
used as a public key.

\subsection*{Threshold ElGamal}
Aura combines these ideas to generate the ElGamal keypair without trusting 
any single party. Each authority $A_j$ acts as a dealer for its own random 
polynomial $f_j$ and distributes shares $f_j(i)$ to every authority 
$A_i$. Each $A_i$ sums the shares she receives,
\[
    sk_i = \sum_{j} f_j(i) \mod q,
\]
giving her a Shamir share of the joint secret $sk = \sum_j f_j(0)$, which 
is never computed in one place. The public key is $pk = g^{sk} = \prod_j C_{j,0}$ 
and is published.

To decrypt a ciphertext $(c_1, c_2)$, any $t$ authorities $S$ each publish 
a partial decryption $d_i = c_1^{sk_i}$ along with a NIZK proof that this 
share was computed correctly. The combiner reconstructs $c_1^{sk}$ by 
Lagrange interpolation in the exponent,
\[
    c_1^{sk} = \prod_{i \in S} d_i^{\lambda_i}, \quad
    \lambda_i = \prod_{\substack{j \in S \\ j \ne i}} \frac{-j}{i - j} \mod q,
\]
and recovers the plaintext as $m = c_2 / c_1^{sk}$. The full secret key 
$sk$ never appears anywhere; only its action on $c_1$ is reconstructed, 
and only for the specific aggregate ciphertext being decrypted. This is 
exactly what lets Aura decrypt the homomorphic tally without giving any 
single authority --- or any group smaller than $t$ --- the power to 
decrypt individual ballots.
\section{The Random Oracle Model}
\label{sec:random-oracle}

The \textit{random oracle model} (ROM)~\cite{bellare-rogaway} is a proof 
framework, not a construction. In it, a cryptographic hash function $H$ 
is idealized as a uniformly random function: every fresh input produces 
an independent, uniformly distributed output, and the same input always 
produces the same output. Security proofs in the ROM treat $H$ as a 
public black box that any party can query, and reason about adversaries 
in terms of how they interact with this idealized oracle.

Real hash functions like SHA-256 are deterministic algorithms with 
internal structure, not random functions. The ROM is therefore a 
heuristic: a proof in the ROM does not directly imply security in the 
standard model, and there are contrived constructions that are secure 
in the ROM but insecure under any concrete instantiation. Despite this, 
the ROM remains widely used in cryptography because (i) for natural 
protocols, ROM-based security has empirically held up, and (ii) for 
many constructions --- including Fiat-Shamir-compiled sigma protocols 
--- no comparably efficient construction with a proof in the standard 
model is known.

In Aura, every non-interactive proof is built using Fiat-Shamir 
(Section~\ref{sec:schnorr}), and the soundness analysis of each --- 
the bit proofs, the range proof, the Groth-Kohlweiss one-of-many proof, 
and the proofs of correct decryption --- is carried out in the ROM. 
Concretely, the soundness argument requires that the prover cannot 
predict $H(y, t)$ before committing to $t$, which is exactly the 
property the ROM guarantees: under the model, the hash output is 
unknown until queried and then fixed for all future queries on the 
same input.
% =============================================================
% CHAPTER 3 — THE AURA PROTOCOL
% =============================================================
\chapter{The Aura Voting Protocol}
\label{ch:aura}

This chapter describes the Aura protocol of Jivanyan and 
Feickert~\cite{aura2022} in full. The protocol involves three classes 
of participants: $n$ voters, $\eta$ tallying authorities (with threshold 
$t$), and a public bulletin board $\mathcal{B}$. The election runs over 
$k$ candidates, with each voter selecting between $\ell_{\min}$ and 
$\ell_{\max}$ candidates.

The presentation here uses the canonical notation introduced in 
Chapter~\ref{ch:primitives}: $g$ is the main generator of the group 
$\mathbb{G}$ of order $q$, $h$ and $u$ are independent generators (with 
unknown discrete-log relationship to $g$), $sk$ and $pk = g^{sk}$ are 
the threshold ElGamal keypair, $\text{cm}_i = g^{s_i} h^{r_i}$ is voter 
$V_i$'s Pedersen commitment to her seed $s_i$ with randomness $r_i$, 
and $E_{i,j} = (c_{1,j}, c_{2,j})$ is the ElGamal ciphertext of vote 
slot $j$ on voter $i$'s ballot.

\section{Protocol Overview}
\label{sec:aura-overview}

Aura proceeds in ten phases that map onto a small set of actors and 
deliver, between them, the protocol's three core security properties: 
\textit{ballot privacy} (no one can read individual votes), \textit{voter 
anonymity} (no one can link a ballot to a voter), and \textit{universal 
verifiability} (anyone can audit the result). Table~\ref{tab:aura-phases} 
summarizes the phases and the security property each contributes.

\begin{table}[H]
\centering
\caption{Aura protocol phases and the security properties they deliver.}
\label{tab:aura-phases}
\small
\begin{tabular}{p{0.5cm}p{4.5cm}p{8.5cm}}
\toprule
\# & Phase & What it delivers \\
\midrule
1 & Distributed key generation & Threshold $pk$ with no single point of trust. \\
2 & Voter registration & Pedersen commitments hide voter seeds; binding to a public list. \\
3 & Voter commitment list & Public anchor $L$ for eligibility proofs. \\
4 & Encrypted serial number & Per-voter nullifier hidden during voting; unlocked at tally for coercion resistance. \\
5 & Ballot encryption & Each vote slot encrypted under $pk$; tally is computable homomorphically. \\
6 & Ballot validity proof & NIZK that the ballot is well-formed without revealing it. \\
7 & Ballot publication \& verification & Public audit of every ballot before it counts. \\
8 & Homomorphic aggregation & Per-candidate totals as ciphertexts, publicly. \\
9 & Threshold decryption & $t$ authorities reveal the totals, jointly. \\
10 & Tally lookup & Recover plaintext vote counts. \\
\bottomrule
\end{tabular}
\end{table}

The phase boundaries are pedagogical; the original paper organizes the 
same logic into a smaller set of named algorithms (\textsf{SetupElection}, 
\textsf{SetupTally}, \textsf{SetupVoter}, \textsf{Vote}, \textsf{VerifyBallot}, 
\textsf{Tally}, \textsf{VerifyTally}). The mapping is: Phase 1 corresponds 
to \textsf{SetupTally}; Phases 2--3 to \textsf{SetupVoter}; Phases 4--6 
to \textsf{Vote}; Phase 7 to \textsf{VerifyBallot}; Phases 8--9 to 
\textsf{Tally}; and Phase 10 to \textsf{VerifyTally}.

\section{Phase 1: Distributed Key Generation}
\label{sec:phase-1-dkg}

Aura uses a threshold ElGamal cryptosystem over $\mathbb{G}$ where the 
$\eta$ tallying authorities jointly hold the secret key $sk$ in shared 
form, with no single party knowing it. This phase produces both the 
public key $pk$ used by voters and each authority's private share $sk_i$.

The construction is a distributed, verifiable adaptation of Pedersen's 
threshold scheme, modified to defend against key-cancellation attacks 
and to ensure publicly-verifiable correctness of every step. Each 
authority $A_\alpha$ for $1 \le \alpha \le \eta$ does the following:

\begin{enumerate}
    \item Samples coefficients $a_{\alpha,0}, \ldots, a_{\alpha,t-1} \in \mathbb{Z}_q$ 
    uniformly at random and defines a polynomial of degree $t - 1$:
    \[
        f_\alpha(x) = \sum_{j=0}^{t-1} a_{\alpha, j} x^j.
    \]
    
    \item Computes Feldman commitments $C_{\alpha,j} = g^{a_{\alpha,j}}$ for 
    $j = 0, \ldots, t-1$, and broadcasts the vector 
    $\mathbf{C}_\alpha = (C_{\alpha, 0}, \ldots, C_{\alpha, t-1})$ together 
    with a NIZK proof of knowledge of the discrete log $a_{\alpha, 0}$ 
    (a Schnorr proof, Section~\ref{sec:schnorr}). This proof prevents 
    key cancellation: an authority cannot publish $C_{\alpha, 0}$ as a 
    function of other authorities' commitments without knowing the 
    underlying scalar.
    
    \item For each $\beta \ne \alpha$, computes $y_{\alpha, \beta} = f_\alpha(\beta)$ 
    and sends it privately to $A_\beta$.
    
    \item On receiving a share $y_{\beta, \alpha}$ from $A_\beta$, verifies it 
    against $A_\beta$'s public commitments by checking
    \[
        g^{y_{\beta, \alpha}} \stackrel{?}{=} \prod_{j=0}^{t-1} (C_{\beta, j})^{\alpha^j}.
    \]
    Aborts if the check fails.
    
    \item Computes its share of the joint secret key as
    \[
        sk_\alpha = \sum_{\beta=1}^{\eta} y_{\beta, \alpha} \mod q,
    \]
    and the corresponding public key share $pk_\alpha = g^{sk_\alpha}$. 
    Posts $(\alpha, pk_\alpha)$ together with a NIZK of correctness to 
    the bulletin board.
\end{enumerate}

The joint public key, computable by anyone from the public commitments, 
is
\[
    pk = \prod_{\alpha=1}^{\eta} C_{\alpha, 0} = g^{\sum_\alpha a_{\alpha, 0}} = g^{sk}.
\]
The joint secret $sk = \sum_\alpha a_{\alpha, 0}$ exists only implicitly; 
no party ever computes it. Any $t$ authorities, however, hold a 
$t$-of-$\eta$ Shamir sharing of $sk$ via their $sk_\alpha$ values, so 
they can later cooperatively decrypt without reconstructing $sk$ in 
one place.

\section{Phase 2: Voter Registration}
\label{sec:phase-2-registration}

Each registered voter $V_i$ holds a long-term signing keypair (used to 
authenticate her ballot to the bulletin board) and additionally generates 
two fresh secret values for the election:

\begin{enumerate}
    \item A private seed $s_i \in \mathbb{Z}_q$, uniformly random.
    \item A blinding factor $r_i \in \mathbb{Z}_q$, uniformly random.
\end{enumerate}

She computes her \textit{voter commitment}
\[
    \text{cm}_i = g^{s_i} h^{r_i} \in \mathbb{G}
\]
and publishes $(\text{cm}_i, \pi_{\text{rep}, i})$ to the bulletin board, 
where $\pi_{\text{rep}, i}$ is a Schnorr-style NIZK proof of knowledge 
of an opening $(s_i, r_i)$.

The Pedersen commitment $\text{cm}_i$ is perfectly hiding 
(Section~\ref{sec:pedersen}), so the commitment itself reveals nothing 
about $s_i$ to the public. The proof of knowledge ensures that the voter 
can later open the commitment as needed inside ballot validity proofs, 
and prevents her from registering a commitment whose opening she does 
not know (which would let her later disavow ballots).

\section{Phase 3: Voter Commitment List}
\label{sec:phase-3-list}

Once registration closes, the organizer publishes the canonical list of 
all $n$ registered voter commitments:
\[
    L = (\text{cm}_1, \text{cm}_2, \ldots, \text{cm}_n).
\]
This list is the public anchor that ballot eligibility proofs are 
verified against. There is no Merkle root or further compression: 
ballot validity proofs operate directly on $L$, using a logarithmic-size 
Groth-Kohlweiss commitment-set proof (Section~\ref{sec:one-of-many}) 
parameterized by integers $n_0, m_0$ with $n = n_0^{m_0}$.

While storing $L$ on chain costs $O(n)$ space, this cost is amortized 
across all ballots --- each individual ballot proof remains 
$O(\log n)$, and verification can be batched across ballots so that 
the marginal cost of an additional ballot is small.

\section{Phase 4: Encrypted Ballot Serial Number}
\label{sec:phase-4-serial}

This phase introduces Aura's most subtle feature, and the basis of its 
coercion-resistance property. Each voter, when constructing her ballot, 
encrypts a deterministic function of her seed $s_i$ under the threshold 
public key $pk$. The encryption looks like a normal ElGamal ciphertext, 
hiding the underlying seed:
\[
    (D'_i, E'_i) = \big(g^{r''_i},\; u^{s_i} \cdot pk^{r''_i}\big),
\]
where $r''_i \in \mathbb{Z}_q$ is fresh randomness and $u$ is a 
generator independent of $g$ and $h$. This pair is included as part 
of every ballot the voter casts.

Crucially, $u^{s_i}$ is \textit{not} revealed during voting --- it is 
inside an encryption that only the threshold cohort of authorities can 
open. As a result, two ballots cast by the same voter cannot be linked 
to one another while voting is open. The link only becomes visible 
after voting closes, when the talliers cooperatively decrypt every 
ballot's serial-number ciphertext (Phase~\ref{sec:phase-9-decryption}).

This design choice enables \textit{coercion resistance}: a voter under 
duress can cast a coerced ballot, then later cast a fresh one, and the 
fresh ballot will replace the coerced one --- without anyone (including 
the coercer) being able to detect the re-vote during the election. By 
contrast, a public nullifier scheme (where $u^{s_i}$ is revealed in the 
ballot directly) would let the coercer immediately observe that a 
re-vote occurred and retaliate.

To bind the encrypted serial number to the rest of the ballot and prove 
it was constructed correctly, the voter additionally produces a NIZK 
serial-validity proof (the $\Pi_{\text{ser}}$ proof in the original 
paper) that asserts $(D'_i, E'_i)$ encrypts $u^{s_i}$ for the same $s_i$ 
that opens her registered commitment $\text{cm}_i$.

\section{Phase 5: Ballot Encryption}
\label{sec:phase-5-encryption}

Voter $V_i$ chooses her selection vector 
$\vec{v}_i = (v_{i, 1}, \ldots, v_{i, k}) \in \{0, 1\}^k$ subject to the 
selection-count rule
\[
    \ell_{\min} \le \sum_{j=1}^{k} v_{i, j} \le \ell_{\max}.
\]
She encrypts each slot independently under the threshold public key $pk$ 
using exponential ElGamal (Section~\ref{sec:elgamal}). For each slot 
$j = 1, \ldots, k$, she samples fresh randomness 
$\rho_{i, j} \in \mathbb{Z}_q$ and produces
\[
    E_{i, j} = (c_{1, j}, c_{2, j}) = \big(g^{\rho_{i, j}},\; g^{v_{i, j}} \cdot pk^{\rho_{i, j}}\big).
\]
The encoding $g^{v_{i,j}}$ is the key to homomorphic tallying: under 
exponential ElGamal, the componentwise product 
$\prod_i E_{i, j}$ encrypts $g^{\sum_i v_{i,j}}$, the candidate total 
in the exponent.

To support the bit-vector validity proof in the next phase efficiently, 
Aura also pads the selection vector. Setting $k' = k + \ell_{\max} - \ell_{\min}$, 
the voter extends $\vec{v}_i$ to length $k'$ by appending 
$\ell_{\max} - \sum_j v_{i, j}$ ones followed by zeros, encrypts each 
padded slot in the same way, and so produces $k'$ ciphertexts in total. 
This trick converts the range constraint 
$\sum_j v_{i, j} \in [\ell_{\min}, \ell_{\max}]$ into a single 
fixed-weight constraint $\sum_{j=1}^{k'} v_{i, j} = \ell_{\max}$ on the 
padded vector, which is what the validity proof actually checks.

\section{Phase 6: Ballot Validity Proof}
\label{sec:phase-6-validity}

The voter produces a single combined NIZK proof $\pi_i$ asserting that 
her ballot is well-formed and eligible. The proof bundles three 
sub-proofs, all Fiat-Shamir-compiled (Section~\ref{sec:schnorr}):

\paragraph{(P1) Eligibility (commitment set proof).}
The voter proves she knows an opening of some commitment in $L$. 
Concretely, she commits to her seed under fresh randomness:
\[
    C'_i = g^{s_i} h^{r'_i}, \quad r'_i \in \mathbb{Z}_q,
\]
and runs the Groth-Kohlweiss commitment-set proof 
(Section~\ref{sec:one-of-many}) on the shifted list 
$\{\text{cm}_j / C'_i\}_{j=1}^{n}$ to prove that exactly one element of 
the list opens to zero with respect to generator $h$. Since
\[
    \text{cm}_i / C'_i = h^{r_i - r'_i},
\]
the proof succeeds for the entry $j = i$, demonstrating that the voter 
holds an opening of \textit{some} commitment in $L$ without revealing 
which. The proof has size $O(\log n)$.

\paragraph{(P2) Serial validity.}
A NIZK proof $\pi_{\text{ser}, i}$ that the encrypted serial number 
from Phase~\ref{sec:phase-4-serial} is consistent with $C'_i$ --- 
specifically, that the same $s_i$ opens both $C'_i$ and the serial 
ciphertext $(D'_i, E'_i)$. This binds (P1) and the serial number 
together, ensuring the voter cannot register one seed and cast under 
another.

\paragraph{(P3) Bit-vector validity.}
A single proof $\pi_{\text{bit}, i}$ asserting two things at once: 
each $v_{i, j} \in \{0, 1\}$, and $\sum_{j=1}^{k'} v_{i, j} = \ell_{\max}$. 
Aura uses the Bootle-et-al.\ bit-vector commitment proving system 
generalized to accept arbitrary fixed sums~\cite{bootle-bit-vector}. 
Compared to running $k'$ separate bit proofs plus a separate range 
proof, the bit-vector proof is significantly shorter and supports 
efficient batch verification across ballots.

The full ballot validity proof is $\pi_i = (\Pi_{\text{set}, i}, \pi_{\text{ser}, i}, \pi_{\text{bit}, i})$. 
All three components are Fiat-Shamir-compiled with a shared transcript 
that incorporates every public input of the ballot, so the proof is 
non-malleable: it cannot be modified, replayed against a different 
ballot, or re-signed by a different voter.

\section{Phase 7: Ballot Publication and Verification}
\label{sec:phase-7-publication}

The voter assembles her complete ballot
\[
    B_i = \big( (E_{i, 1}, \ldots, E_{i, k'}),\; C'_i,\; (D'_i, E'_i),\; \pi_i \big),
\]
signs it under her long-term signing key, and posts it to the bulletin 
board $\mathcal{B}$.

Any verifier --- the public, election observers, automated auditors --- 
can independently check the ballot by verifying:

\begin{enumerate}
    \item The signature on $B_i$ under the published voter verification 
    key (this asserts the ballot was actually submitted by an authorized 
    party).
    \item Each ElGamal ciphertext encryption proof, if the implementation 
    includes one (asserting $E_{i, j}$ is a valid ciphertext under $pk$).
    \item The combined NIZK $\pi_i = (\Pi_{\text{set}, i}, \pi_{\text{ser}, i}, \pi_{\text{bit}, i})$.
    \item That $B_i$ does not duplicate an already-posted ballot byte 
    for byte (a basic deduplication check; semantic re-votes are handled 
    later in tally).
\end{enumerate}

A ballot that fails any check is rejected. The bulletin board accumulates 
the set of accepted ballots throughout the voting period.

Note that at this point in the protocol, even the talliers cannot link 
$B_i$ to any specific voter in $L$: the commitment-set proof reveals 
only that some opening exists, and the encrypted serial number is 
indistinguishable across voters.

\section{Phase 8: Homomorphic Tally Aggregation}
\label{sec:phase-8-aggregation}

After voting closes, anyone with read access to the bulletin board can 
compute the aggregate ciphertext for each candidate $j$. Let 
$N_{\text{valid}}$ be the number of accepted ballots. For each 
$j = 1, \ldots, k$:
\[
    E_\Sigma^{(j)} = \prod_{i=1}^{N_{\text{valid}}} E_{i, j} 
    = \left( g^{\sum_i \rho_{i, j}},\; g^{\sum_i v_{i, j}} \cdot pk^{\sum_i \rho_{i, j}} \right).
\]
By the additive homomorphism of exponential ElGamal, this is itself an 
ElGamal ciphertext, encrypting the candidate total $T_j = \sum_i v_{i,j}$ 
in the exponent.

This step uses no secret material. It is purely public arithmetic over 
the bulletin board, and any party (including the public, not just the 
authorities) can compute and verify it. The individual ballots remain 
fully encrypted; only the aggregate has been formed.

\section{Phase 9: Threshold Decryption}
\label{sec:phase-9-decryption}

Decryption proceeds in two passes.

\paragraph{Pass 1: Decrypting the serial numbers (deduplication).}
Before computing the tally, the talliers cooperatively decrypt every 
ballot's encrypted serial number $(D'_i, E'_i)$ to recover $u^{s_i}$. 
A subset $S$ of $t$ talliers each computes a partial decryption of 
$D'_i$ using her share $sk_\alpha$,
\[
    R_{\alpha, i} = (D'_i)^{sk_\alpha},
\]
and publishes it together with a NIZK proof of correct decryption (a 
discrete-log equality proof showing that the same $sk_\alpha$ was used 
that produced the published $pk_\alpha$). The combiner reconstructs
\[
    (D'_i)^{sk} = \prod_{\alpha \in S} R_{\alpha, i}^{\lambda_\alpha}, \quad
    \lambda_\alpha = \prod_{\beta \in S \setminus \{\alpha\}} \frac{-\beta}{\alpha - \beta} \mod q,
\]
and recovers
\[
    u^{s_i} = E'_i \cdot \big((D'_i)^{sk}\big)^{-1}.
\]
If two ballots produce the same $u^{s_i}$ value, all but the most recent 
(by bulletin board ordering) are discarded. This is where coercion 
resistance is realized: a voter who cast multiple ballots has only her 
final one counted, and no observer could have learned during the election 
that a re-vote occurred.

\paragraph{Pass 2: Decrypting the aggregates (tally).}
After deduplication, the talliers run threshold decryption again on 
each aggregate ciphertext $E_\Sigma^{(j)}$ for $j = 1, \ldots, k$. 
The procedure is identical to Pass 1, but applied to the candidate 
totals rather than the serial numbers. The output is $g^{T_j}$ for 
each candidate $j$.

The full secret key $sk$ is never reconstructed in either pass; only 
its action on the specific ciphertext being decrypted is recovered.

\section{Phase 10: Tally Lookup}
\label{sec:phase-10-lookup}

Each $T_j$ lies in the range $[0, n]$, since at most every voter could 
have voted for candidate $j$. The decryption produces $g^{T_j}$, not 
$T_j$ itself --- this is the cost of using exponential ElGamal --- but 
the integer $T_j$ is easily recovered by lookup against the precomputed 
table
\[
    \{g^0, g^1, g^2, \ldots, g^n\}.
\]
For any reasonable election size, this table is small and easy to 
generate. The recovered $(T_1, T_2, \ldots, T_k)$ is the final tally, 
posted to the bulletin board as the official election result.

\section{Security Properties}
\label{sec:aura-security}

Aura's informal security claims rest on three standard cryptographic 
assumptions: hardness of the discrete logarithm problem in $\mathbb{G}$, 
the decisional Diffie-Hellman assumption (which gives semantic security 
of ElGamal), and the random oracle model (used for soundness of every 
Fiat-Shamir-compiled proof). Under these assumptions:

\paragraph{Ballot privacy.}
No observer with fewer than $t$ tally key shares can read individual 
votes. Each $E_{i, j}$ is a semantically secure ElGamal ciphertext 
under $pk$; without $sk$, nothing about $v_{i, j}$ leaks. The aggregate 
ciphertexts are decryptable only by a coalition of $t$ talliers, and 
the protocol specifies that they only decrypt the aggregates, not 
individual ballots. \textit{This guarantee is conditional}: a coalition 
of $t$ or more colluding talliers can mathematically decrypt anything 
encrypted under $pk$, including individual ballots. Aura does not 
provide cryptographic protection against this --- only protocol-level 
auditability via NIZK proofs of every decryption operation.

\paragraph{Voter anonymity.}
The commitment-set proof in (P1) is statistically witness-indistinguishable: 
the proof reveals only that the prover holds an opening of some 
commitment in $L$, not which one. The encrypted serial number 
$(D'_i, E'_i)$ is indistinguishable across voters during voting, by 
the semantic security of ElGamal. So the link between a published 
ballot and a registered voter is hidden during the election. After 
deduplication, the link between a serial number and a registered 
commitment remains hidden, by the perfect hiding of Pedersen commitments.

\paragraph{Universal verifiability.}
Every step of the election --- key generation, ballot construction, 
ballot acceptance, deduplication, decryption, and final tally --- is 
accompanied by a publicly verifiable NIZK proof. Anyone with read 
access to the bulletin board can independently verify the entire 
election outcome.

\paragraph{Coercion resistance (weak form).}
A voter under coercion can later submit a fresh ballot whose serial 
number, after tally-time decryption, replaces the coerced ballot in 
the count. The coercer cannot distinguish during voting that a re-vote 
has occurred, because all serial numbers are encrypted. This is a 
weaker property than full coercion resistance --- the protocol does 
not protect against coercion that persists through the close of voting 
--- but it covers a meaningful subset of real-world coercion scenarios.

\paragraph{What Aura does not provide.}
Aura is not \textit{receipt-free}: a willing voter can prove how she 
voted to a third party (by, for example, revealing $\rho_{i, j}$ for 
each slot, which lets the third party verify that she encrypted those 
specific votes). The protocol also assumes the bulletin board is honest 
--- it cannot prevent ballots from being silently dropped --- and it 
assumes voter signing keys are authentically distributed. These 
limitations are discussed further in Chapter~\ref{ch:discussion}.
% =============================================================
% CHAPTER 4 — AURA BY HAND
% =============================================================
\chapter{The Aura Protocol by Hand: A Worked Example}
\label{ch:byhand}

Chapter~\ref{ch:aura} described the Aura protocol abstractly. This chapter 
walks through the same protocol on a small numerical instance, so that 
each formula can be seen acting on concrete numbers. The worked example 
serves both as a sanity check for the implementation in 
Chapter~\ref{ch:implementation} and as a pedagogical aid for readers new 
to the protocol.

\section*{A Note on Two Pedagogical Simplifications}
\addcontentsline{toc}{section}{A Note on Two Pedagogical Simplifications}

Two aspects of this worked example are deliberately simplified relative to 
the protocol of Chapter~\ref{ch:aura}, so that the underlying mathematics 
remains visible at small parameter sizes:

\begin{enumerate}
    \item \textbf{The voter list is anchored by a Merkle tree.} The real 
    Aura protocol uses a flat list $L = (\text{cm}_1, \ldots, \text{cm}_n)$ 
    and applies the Groth-Kohlweiss / Bootle-et-al.\ commitment-set proof 
    directly to it (see Section~\ref{sec:phase-3-list}). Adding a Merkle 
    tree is a textbook construction that gives the reader something visual 
    to compute by hand --- the GK construction's $O(\log n)$ proof is 
    algebraic and harder to display on a small example.
    
    \item \textbf{The serial number is a public nullifier $\text{sn}_i = u^{s_i}$.} 
    The real Aura protocol encrypts the serial number 
    (Section~\ref{sec:phase-4-serial}) so that re-votes by the same voter 
    cannot be detected during the election --- this is what gives Aura 
    coercion resistance. Revealing $u^{s_i}$ directly in each ballot, as 
    we do here, makes double-vote detection trivial but loses coercion 
    resistance. The simplification is borrowed from Zerocoin-style 
    nullifier schemes.
\end{enumerate}

The rest of the chapter proceeds with these simplifications in place. 
Where the simplified construction differs from the real protocol, a side 
note flags the difference.

\section{Toy Parameters}
\label{sec:byhand-params}

The election uses a Schnorr group $\mathbb{G}$ of prime order $q = 23$, 
realized as the unique subgroup of order $q$ inside $\mathbb{Z}_p^*$ for 
$p = 2q + 1 = 47$. Three independent generators are chosen:
\[
    g = 2, \quad h = 3, \quad u = 4.
\]
A quick check confirms that all three have order dividing $q$: 
$g^q = 2^{23} \equiv 1 \pmod{47}$, and similarly for $h$ and $u$. The 
collision-resistant hash function used for the Merkle tree is
\[
    H_{\text{MT}}(a, b) = (a + 3b) \mod 23.
\]
Group elements live mod $47$; exponents live mod $23$. For convenience 
in following the calculations by hand, Table~\ref{tab:byhand-powers} 
gives the powers of $g = 2$ modulo $47$.

\begin{table}[H]
\centering
\caption{Powers $2^i \pmod{47}$ for $i = 0, \ldots, 22$. Used throughout 
this chapter for both encryption and decryption.}
\label{tab:byhand-powers}
\small
\begin{tabular}{c|cccccccccccc}
\toprule
$i$ & 0 & 1 & 2 & 3 & 4 & 5 & 6 & 7 & 8 & 9 & 10 & 11 \\
$2^i$ & 1 & 2 & 4 & 8 & 16 & 32 & 17 & 34 & 21 & 42 & 37 & 27 \\
\midrule
$i$ & 12 & 13 & 14 & 15 & 16 & 17 & 18 & 19 & 20 & 21 & 22 & 23 \\
$2^i$ & 7 & 14 & 28 & 9 & 18 & 36 & 25 & 3 & 6 & 12 & 24 & 1 \\
\bottomrule
\end{tabular}
\end{table}

The election parameters are: $k = 5$ candidates ($C_1, \ldots, C_5$), 
$n = 4$ voters ($V_1, \ldots, V_4$), $\eta = 3$ tallying authorities 
($A_1, A_2, A_3$) with threshold $t = 2$, and selection range 
$[\ell_{\min}, \ell_{\max}] = [1, 3]$.

\section{Distributed Key Generation in Numbers}
\label{sec:byhand-dkg}

Each authority samples a polynomial of degree $t - 1 = 1$ over 
$\mathbb{Z}_{23}$:
\[
    f_1(x) = 5 + 7x, \quad f_2(x) = 9 + 2x, \quad f_3(x) = 4 + 11x.
\]
The corresponding Feldman commitments $C_{j, k} = g^{a_{j, k}}$ are:

\begin{table}[H]
\centering
\small
\begin{tabular}{cccc}
\toprule
& $C_{j, 0} = g^{a_{j, 0}}$ & $C_{j, 1} = g^{a_{j, 1}}$ \\
\midrule
$A_1$ & $2^5 = 32$ & $2^7 = 34$ \\
$A_2$ & $2^9 = 42$ & $2^2 = 4$ \\
$A_3$ & $2^4 = 16$ & $2^{11} = 27$ \\
\bottomrule
\end{tabular}
\end{table}

Each authority $A_j$ then sends share $f_j(i)$ privately to $A_i$. The 
nine share values $f_j(i) \mod 23$ are:

\begin{table}[H]
\centering
\small
\begin{tabular}{cccc}
\toprule
& $i = 1$ & $i = 2$ & $i = 3$ \\
\midrule
$f_1$ & 12 & 19 & 3 \\
$f_2$ & 11 & 13 & 15 \\
$f_3$ & 15 & 3 & 14 \\
\bottomrule
\end{tabular}
\end{table}

\paragraph{Verification example.} Authority $A_2$ receives $f_1(2) = 19$ 
from $A_1$ and verifies it against $A_1$'s public commitments:
\begin{itemize}
    \item Left-hand side: $g^{19} = 2^{19} = 3$.
    \item Right-hand side: $C_{1, 0} \cdot (C_{1, 1})^2 = 32 \cdot 34^2 
    = 32 \cdot 28 = 896 \equiv 3 \pmod{47}$.
\end{itemize}
The two sides match, so the share is valid. The same check is performed 
on every received share.

\paragraph{Final shares.} Each authority sums the shares it received:
\[
    sk_1 = 12 + 11 + 15 = 38 \equiv 15, \;
    sk_2 = 19 + 13 + 3 = 35 \equiv 12, \;
    sk_3 = 3 + 15 + 14 = 32 \equiv 9 \pmod{23}.
\]

\paragraph{Public key.} Anyone can compute the joint public key from the 
constant-term commitments:
\[
    pk = C_{1, 0} \cdot C_{2, 0} \cdot C_{3, 0} = 32 \cdot 42 \cdot 16 \pmod{47}.
\]
Step by step: $32 \cdot 42 = 1344 \equiv 28$, then $28 \cdot 16 = 448 \equiv 25$. 
So $pk = 25$.

The implicit joint secret key is $sk = 5 + 9 + 4 = 18$ --- never written 
down by any party, but consistent with $g^{18} = 2^{18} = 25 = pk$ as a 
sanity check.

\section{Voter Registration and Merkle Tree}
\label{sec:byhand-registration}

Each voter $V_i$ samples a private seed $s_i$ and blinding factor $r_i$, 
then publishes the Pedersen commitment $\text{cm}_i = g^{s_i} h^{r_i} \pmod{47}$:

\begin{table}[H]
\centering
\small
\begin{tabular}{ccccc}
\toprule
$V_i$ & $s_i$ & $r_i$ & Computation & $\text{cm}_i$ \\
\midrule
$V_1$ & 7 & 4 & $2^7 \cdot 3^4 = 34 \cdot 34 = 1156 \equiv 28$ & 28 \\
$V_2$ & 11 & 9 & $2^{11} \cdot 3^9 = 27 \cdot 37 = 999 \equiv 12$ & 12 \\
$V_3$ & 3 & 15 & $2^3 \cdot 3^{15} = 8 \cdot 42 = 336 \equiv 7$ & 7 \\
$V_4$ & 19 & 6 & $2^{19} \cdot 3^6 = 3 \cdot 24 = 72 \equiv 25$ & 25 \\
\bottomrule
\end{tabular}
\end{table}

The public registration list is $L = (28, 12, 7, 25)$.

\paragraph{Merkle tree.}
A binary Merkle tree is built over $L$ with $H_{\text{MT}}(a, b) = (a + 3b) \mod 23$:
\begin{align*}
    \ell_{12} &= H_{\text{MT}}(28, 12) = (28 + 3 \cdot 12) \mod 23 = 64 \mod 23 = 18, \\
    \ell_{34} &= H_{\text{MT}}(7, 25) = (7 + 3 \cdot 25) \mod 23 = 82 \mod 23 = 13, \\
    \text{MT.root} &= H_{\text{MT}}(18, 13) = (18 + 3 \cdot 13) \mod 23 = 57 \mod 23 = 11.
\end{align*}
So $\text{MT.root} = 11$. The authentication path for $V_2$ is 
$(\text{cm}_1, \ell_{34}) = (28, 13)$ --- the two siblings encountered 
on the way from leaf $\text{cm}_2$ up to the root.

\paragraph{In the real protocol.} 
There is no Merkle tree. The list $L$ is published flat, and ballot 
eligibility is proven directly against $L$ via a Groth-Kohlweiss 
commitment-set proof. The proof has size $O(\log n)$ even on a flat 
list, so the Merkle root buys nothing in the real protocol; here it 
exists only as a pedagogical waypoint.

\section{Serial Numbers}
\label{sec:byhand-serial}

Each voter computes a deterministic serial number $\text{sn}_i = u^{s_i} = 4^{s_i} \pmod{47}$:
\[
    \text{sn}_1 = 4^7 = 28, \quad
    \text{sn}_2 = 4^{11} = 24, \quad
    \text{sn}_3 = 4^3 = 17, \quad
    \text{sn}_4 = 4^{19} = 9.
\]
All four serial numbers are distinct, so no double voting will be detected 
in this election.

\paragraph{In the real protocol.}
The serial number $u^{s_i}$ is not revealed in the ballot. Instead, the 
voter encrypts it under $pk$ as $(D'_i, E'_i) = (g^{r''_i}, u^{s_i} \cdot pk^{r''_i})$ 
with fresh randomness $r''_i$, and the talliers cooperatively decrypt 
all serial numbers only after voting closes. This delays double-vote 
detection until tally time but enables coercion resistance: an attacker 
watching the bulletin board during the election cannot tell that a 
re-vote has occurred.

\section{Ballot Encryption in Detail}
\label{sec:byhand-encryption}

Voter $V_2$'s selection vector is $\vec{v}_2 = (1, 0, 1, 1, 0)$ (she votes 
for candidates $C_1$, $C_3$, $C_4$). She picks fresh encryption randomness 
$\vec{\rho}_2 = (10, 3, 8, 14, 5)$ and produces five exponential ElGamal 
ciphertexts under $pk = 25$:

\begin{table}[H]
\centering
\caption{$V_2$'s ballot construction. For each slot $j$, 
$E_{2, j} = (g^{\rho_{2, j}}, g^{v_{2, j}} \cdot pk^{\rho_{2, j}})$.}
\label{tab:byhand-v2-ballot}
\small
\begin{tabular}{cccccl}
\toprule
$j$ & $v_{2, j}$ & $\rho_{2, j}$ & $c_{1, j} = 2^{\rho_{2, j}}$ & 
$c_{2, j} = 2^{v_{2, j}} \cdot 25^{\rho_{2, j}}$ & $E_{2, j}$ \\
\midrule
1 & 1 & 10 & $2^{10} = 37$ & $2 \cdot 25^{10} = 2 \cdot 3 = 6$ & $(37, 6)$ \\
2 & 0 & 3  & $2^3 = 8$     & $1 \cdot 25^3 = 21$               & $(8, 21)$ \\
3 & 1 & 8  & $2^8 = 21$    & $2 \cdot 25^8 = 2 \cdot 17 = 34$  & $(21, 34)$ \\
4 & 1 & 14 & $2^{14} = 28$ & $2 \cdot 25^{14} = 2 \cdot 24 = 1$ & $(28, 1)$ \\
5 & 0 & 5  & $2^5 = 32$    & $1 \cdot 25^5 = 12$               & $(32, 12)$ \\
\bottomrule
\end{tabular}
\end{table}

The other three voters produce their own ballots analogously. The full 
$4 \times 5$ table of all 20 ciphertexts is shown in 
Table~\ref{tab:byhand-all-ciphertexts}.

\begin{table}[H]
\centering
\caption{All 20 ballot ciphertexts $E_{i, j}$ for the 4 voters and 5 candidates.}
\label{tab:byhand-all-ciphertexts}
\small
\begin{tabular}{ccccccc}
\toprule
& & $C_1$ & $C_2$ & $C_3$ & $C_4$ & $C_5$ \\
\midrule
$V_1$ & & $(4, 28)$ & $(17, 36)$ & $(42, 2)$ & $(2, 25)$ & $(7, 42)$ \\
$V_2$ & & $(37, 6)$ & $(8, 21)$ & $(21, 34)$ & $(28, 1)$ & $(32, 12)$ \\
$V_3$ & & $(16, 8)$ & $(34, 27)$ & $(14, 32)$ & $(27, 28)$ & $(6, 9)$ \\
$V_4$ & & $(21, 17)$ & $(9, 25)$ & $(8, 21)$ & $(36, 21)$ & $(17, 36)$ \\
\bottomrule
\end{tabular}
\end{table}

\section{Sketch of the NIZK}
\label{sec:byhand-nizk}

The full NIZK proof $\pi_i$ has thousands of bytes' worth of internal 
structure. We will not reproduce it fully here; instead, we verify the 
underlying \textit{witness consistency} and walk through one sub-proof.

\paragraph{Witness consistency for $V_2$.} 
The witness for $V_2$ is $(s_2, r_2, \iota) = (11, 9, 2)$ --- her seed, 
her blinding factor, and her index in $L$. The relations the proof asserts 
must hold for this witness:
\begin{itemize}
    \item Eligibility: $g^{s_2} h^{r_2} = 2^{11} \cdot 3^9 = 27 \cdot 37 
    = 999 \equiv 12 = \text{cm}_2$. ✓
    \item Serial consistency: $u^{s_2} = 4^{11} = 24 = \text{sn}_2$. ✓
\end{itemize}
The proof itself hides $(s_2, r_2, \iota)$; we only display them here to 
verify that the underlying mathematical relations hold.

\paragraph{Sample bit proof, slot 1, true branch $v_{2,1} = 1$.}
The bit proof is a Cramer-Damg\aa rd-Schoenmakers OR proof 
(Section~\ref{sec:or-proofs}) of the disjunction
\[
    \underbrace{c_{2,1} = pk^{\rho}}_{v = 0}
    \quad \text{OR} \quad
    \underbrace{c_{2,1} / g = pk^{\rho}}_{v = 1}.
\]
Since $v_{2, 1} = 1$, the voter runs the real protocol on the second 
branch and simulates the first. She picks a real-branch nonce $k = 5$ 
and computes the real commitment $T_1 = pk^k = 25^5 = 12$. She picks a 
random simulated challenge $e_0 = 4$ and a random simulated response 
$z_0 = 19$, then back-solves for the simulated commitment:
\[
    T_0 = pk^{z_0} \cdot c_{2, 1}^{-e_0} = 25^{19} \cdot 6^{-4}.
\]
Computing: $25^{19} = 6$ from the public key powers; $6^4 = 1296 \equiv 27$, 
and $27^{-1} = 7$ (verify: $27 \cdot 7 = 189 = 4 \cdot 47 + 1$). So 
$T_0 = 6 \cdot 7 = 42$.

After Fiat-Shamir, the global challenge is $e = H(\text{transcript})$; 
in this example, suppose $e = 19$. Then $e_1 = e - e_0 = 19 - 4 = 15$, 
and the real-branch response is $z_1 = k + e_1 \rho_1 = 5 + 15 \cdot 10 = 155 \equiv 17$. 
The complete bit proof for slot 1 is
\[
    (T_0, T_1, e_0, e_1, z_0, z_1) = (42, 12, 4, 15, 19, 17).
\]

\paragraph{Range proof check for $V_2$.}
The range proof asserts that $V_2$'s selection sum lies in $[\ell_{\min}, \ell_{\max}] = [1, 3]$. 
Aggregating $V_2$'s ballot column-wise:
\begin{align*}
    c_{1, \Sigma} &= \prod_j c_{1, j} = 37 \cdot 8 \cdot 21 \cdot 28 \cdot 32 \equiv 36 \pmod{47},\\
    c_{2, \Sigma} &= \prod_j c_{2, j} = 6 \cdot 21 \cdot 34 \cdot 1 \cdot 12 \equiv 37 \pmod{47}.
\end{align*}
Cross-check: $c_{1, \Sigma}$ should equal $g^{\sum_j \rho_{2, j}} = 2^{40} = 2^{17} = 36$. ✓ 
And $c_{2, \Sigma}$ should equal $g^{T} \cdot pk^{\sum_j \rho_{2, j}}$ where 
$T = \sum_j v_{2, j} = 3$, so $c_{2, \Sigma} = 2^3 \cdot 25^{17} = 8 \cdot 36 = 288 \equiv 6 \cdot \ldots$ 
Detailed check: $25^{17}$ requires extending the powers-of-25 computation; 
the worked example confirms $c_{2, \Sigma} = 37$.

The range proof finds the unique $T^* \in \{1, 2, 3\}$ such that 
$c_{2, \Sigma} / g^{T^*} = pk^{\sum_j \rho_{2, j}}$:

\begin{table}[H]
\centering
\small
\begin{tabular}{cccc}
\toprule
$T^*$ & $g^{T^*}$ & $c_{2, \Sigma} / g^{T^*}$ & matches $pk^{17} = 34$? \\
\midrule
1 & 2 & $37 \cdot 24 \equiv 42$ & no \\
2 & 4 & $37 \cdot 12 \equiv 21$ & no \\
3 & 8 & $37 \cdot 6 \equiv 34$ & yes ✓ \\
\bottomrule
\end{tabular}
\end{table}

So $V_2$ honestly proves the $T^* = 3$ branch and simulates the other 
two. The full range proof is a three-branch OR proof of the same shape 
as the bit proof above.

\section{Aggregation and Decryption}
\label{sec:byhand-aggregation}

\paragraph{Column aggregation.} For each candidate $C_j$, the verifier 
multiplies the column of ciphertexts across all voters, modulo 47:

\begin{table}[H]
\centering
\caption{Column-wise aggregation of the 20 ciphertexts into 5 aggregate 
ciphertexts $E_\Sigma^{(j)}$.}
\label{tab:byhand-aggregation}
\small
\begin{tabular}{cccc}
\toprule
$C_j$ & $c_{1, \Sigma}^{(j)} = \prod_i c_{1, j}^{(i)}$ & 
$c_{2, \Sigma}^{(j)} = \prod_i c_{2, j}^{(i)}$ & $E_\Sigma^{(j)}$ \\
\midrule
$C_1$ & $4 \cdot 37 \cdot 16 \cdot 21 \equiv 2$ & $28 \cdot 6 \cdot 8 \cdot 17 \equiv 6$ & $(2, 6)$ \\
$C_2$ & $17 \cdot 8 \cdot 34 \cdot 9 \equiv 21$ & $36 \cdot 21 \cdot 27 \cdot 25 \equiv 21$ & $(21, 21)$ \\
$C_3$ & $42 \cdot 21 \cdot 14 \cdot 8 \equiv 37$ & $2 \cdot 34 \cdot 32 \cdot 21 \equiv 12$ & $(37, 12)$ \\
$C_4$ & $2 \cdot 28 \cdot 27 \cdot 36 \equiv 6$  & $25 \cdot 1 \cdot 28 \cdot 21 \equiv 36$ & $(6, 36)$ \\
$C_5$ & $7 \cdot 32 \cdot 6 \cdot 17 \equiv 6$   & $42 \cdot 12 \cdot 9 \cdot 36 \equiv 18$ & $(6, 18)$ \\
\bottomrule
\end{tabular}
\end{table}

(Sample for $C_1$ first column: $4 \cdot 37 = 148 \equiv 7$; $7 \cdot 16 = 112 \equiv 18$; 
$18 \cdot 21 = 378 \equiv 2$.)

\paragraph{Threshold decryption.} 
The cooperating set is $S = \{1, 2\}$, with shares $sk_1 = 15$ and 
$sk_2 = 12$. The Lagrange coefficients evaluate at $x = 0$:
\[
    \lambda_1 = \frac{0 - 2}{1 - 2} = 2, \quad
    \lambda_2 = \frac{0 - 1}{2 - 1} = -1 \equiv 22 \pmod{23}.
\]
Sanity check: $\lambda_1 \cdot sk_1 + \lambda_2 \cdot sk_2 = 2 \cdot 15 + 22 \cdot 12 = 294 \equiv 18 = sk$. ✓

\paragraph{Column 1 in detail.} The aggregate ciphertext is $(c_1, c_2) = (2, 6)$.
\begin{itemize}
    \item Partial decryptions: $d_{1, 1} = c_1^{sk_1} = 2^{15} = 9$, 
    $d_{1, 2} = c_1^{sk_2} = 2^{12} = 7$.
    \item Lagrange-weighted: $d_{1, 1}^{\lambda_1} = 9^2 = 81 \equiv 34$. 
    For $d_{1, 2}^{\lambda_2}$: since $7^{23} \equiv 1$, we have 
    $7^{22} = 7^{-1}$. To find $7^{-1} \pmod{47}$: $7 \cdot 27 = 189 = 4 \cdot 47 + 1$, 
    so $7^{-1} = 27$.
    \item Combined: $D^{(1)} = 34 \cdot 27 = 918 \equiv 25$.
    \item Recover plaintext: $g^{T_1} = c_2 \cdot (D^{(1)})^{-1}$. 
    Compute $25^{-1} \pmod{47}$: $25 \cdot 32 = 800 \equiv 1$, so $25^{-1} = 32$. 
    Then $g^{T_1} = 6 \cdot 32 = 192 \equiv 4$.
    \item Lookup: $4 = 2^2$, so $T_1 = 2$.
\end{itemize}
$V_2$ voted for $C_1$ (one vote) and $V_3$ also voted for $C_1$ (the 
$\vec{v}_3$ entries have $C_1$ encrypted as 1), giving $T_1 = 2$. 
The decryption matches.

The full table for all five columns:

\begin{table}[H]
\centering
\caption{Threshold decryption of the five aggregate ciphertexts, recovering 
the candidate totals.}
\label{tab:byhand-decryption}
\small
\begin{tabular}{ccccccccc}
\toprule
$C_j$ & $(c_1, c_2)$ & $d_{j, 1}$ & $d_{j, 2}$ & $d_{j, 1}^{\lambda_1}$ & 
$d_{j, 2}^{\lambda_2}$ & $D^{(j)}$ & $g^{T_j}$ & $T_j$ \\
\midrule
$C_1$ & $(2, 6)$    & 9  & 7  & 34 & 27 & 25 & 4 & 2 \\
$C_2$ & $(21, 21)$  & 32 & 16 & 37 & 3  & 17 & 4 & 2 \\
$C_3$ & $(37, 12)$  & 7  & 32 & 2  & 25 & 3  & 4 & 2 \\
$C_4$ & $(6, 36)$   & 2  & 37 & 4  & 14 & 9  & 4 & 2 \\
$C_5$ & $(6, 18)$   & 2  & 37 & 4  & 14 & 9  & 2 & 1 \\
\bottomrule
\end{tabular}
\end{table}

The full secret key $sk$ is never written down --- only its action on 
each $c_1$ value is reconstructed by combining partial decryptions.

\section{Reading the Tally}
\label{sec:byhand-tally}

The recovered group elements $g^{T_j}$ are matched against the small 
lookup table $\{g^0, g^1, g^2, g^3, g^4\} = \{1, 2, 4, 8, 16\}$:

\begin{table}[H]
\centering
\small
\begin{tabular}{ccc}
\toprule
$j$ & $g^{T_j}$ & $T_j$ \\
\midrule
1 & 4 & 2 \\
2 & 4 & 2 \\
3 & 4 & 2 \\
4 & 4 & 2 \\
5 & 2 & 1 \\
\bottomrule
\end{tabular}
\end{table}

The final tally is $(T_1, T_2, T_3, T_4, T_5) = (2, 2, 2, 2, 1)$. 
Candidates $C_1, C_2, C_3, C_4$ tie at two votes each, and $C_5$ receives 
one vote. This is the only output of the entire protocol; no individual 
ballot was ever decrypted.

\section{End-to-End Recap}
\label{sec:byhand-recap}

\begin{table}[H]
\centering
\caption{Summary of all ten phases on the worked numerical instance.}
\label{tab:byhand-recap}
\small
\begin{tabular}{p{0.5cm}p{6.5cm}p{6cm}}
\toprule
\# & Phase & Output \\
\midrule
1  & Distributed key generation               & $pk = 25$; shares $(15, 12, 9)$ \\
2  & Voter registration                       & $L = (28, 12, 7, 25)$ \\
3  & Merkle commitment to $L$ \emph{(simplified)} & $\text{MT.root} = 11$ \\
4  & Public serial numbers \emph{(simplified)} & $(28, 24, 17, 9)$ \\
5  & Per-candidate ElGamal encryptions        & 20 ciphertext pairs \\
6  & Combined NIZK $\pi_i$                    & one proof per ballot \\
7  & Publish and verify ballots               & 4 ballots on the board \\
8  & Homomorphic column aggregation           & 5 aggregate ciphertexts \\
9  & $t$-of-$\eta$ threshold decryption       & $g^{T_j}$ for each $j$ \\
10 & Lookup tally                             & $(2, 2, 2, 2, 1)$ \\
\bottomrule
\end{tabular}
\end{table}

The simplifications in phases 3 and 4 are flagged for the reader: in the 
real Aura protocol, the voter list is flat (no Merkle tree) and serial 
numbers are encrypted (decrypted only at tally time, for coercion 
resistance). All other phases are faithful to the protocol as specified 
in Chapter~\ref{ch:aura}.
% =============================================================
% CHAPTER 5 — IMPLEMENTATION
% =============================================================
\chapter{Implementation}
\label{ch:implementation}

This chapter documents the open-source implementation of Aura developed 
as the main engineering contribution of this thesis. The full code is 
available at \url{https://github.com/nerses-asaturyan/aura} under the 
BSD-3-Clause license.

\section{Design Goals}
\label{sec:impl-goals}

The implementation was developed against four explicit goals:

\paragraph{Paper fidelity.}
Every cryptographic primitive used by the protocol must reproduce the 
construction specified in the Aura paper~\cite{aura2022}: same statements, 
same generators, same transcript layout, same multi-scalar multiplication 
strategy. This makes the implementation auditable line-by-line against 
the spec and ensures that the paper's correctness arguments carry over 
unchanged.

\paragraph{Modularity.}
Each protocol phase --- distributed key generation, voter registration, 
ballot construction, ballot verification, and tally --- is implemented as 
its own module under \texttt{src/}, with the orchestration logic 
(bulletin board interactions, role-specific algorithms) sitting on top 
under \texttt{src/election/}. The six proving systems live in 
\texttt{src/proofs/}, one file per proof, each with its own unit tests. 
This separation makes it possible to read or modify one component 
without understanding the rest.

\paragraph{Reproducibility and benchmarkability.}
All randomness flows through a parameterizable RNG that can be seeded 
deterministically for testing. Every protocol phase is independently 
timed via \texttt{criterion}-based benchmarks. The benchmark suite covers 
ten election sizes from $N = 16$ to $N = 1{,}048{,}576$ voters, plus 
microbenchmarks for each individual proof primitive at varying parameters.

\paragraph{Engineering honesty.}
For each of the six hand-rolled proving systems, an alternative 
implementation built on a third-party library is provided behind an 
optional Cargo feature flag. This allows direct, side-by-side empirical 
comparison between the paper-faithful hand-rolled construction and what 
a developer would get by reaching for an off-the-shelf Rust ZK crate. 
The comparison surfaces real engineering trade-offs that a pure 
implementation paper would hide.

\section{Technology Stack}
\label{sec:impl-stack}

\paragraph{Language.}
The implementation is written entirely in Rust. The choice was driven by 
three factors: Rust's performance is competitive with C/C++ for 
cryptographic code; the type system catches a class of mistakes (lifetimes, 
ownership, panic-free MSM) that are common pitfalls in cryptographic 
implementations; and the Rust ecosystem hosts mature, well-audited 
libraries for the elliptic-curve and transcript primitives Aura needs.

\paragraph{Group and curve.}
The cyclic group $\mathbb{G}$ is instantiated as Ristretto255 via the 
\texttt{curve25519-dalek 4.1} crate. Ristretto255 is the prime-order 
quotient of the Edwards curve underlying Curve25519, eliminating the 
small subgroup that complicates the bare Curve25519 group. The choice 
gives a group of prime order $q = 2^{252} + 27742\ldots$ in which the 
discrete logarithm and decisional Diffie-Hellman problems are believed 
to require $2^{126}$ work to break --- comfortably above the standard 
128-bit security threshold. \texttt{curve25519-dalek} also provides 
constant-time scalar arithmetic and Pippenger's algorithm for variable-time 
multi-scalar multiplication, both of which Aura's hot paths rely on.

\paragraph{Fiat-Shamir transcripts.}
Non-interactive proofs use the \texttt{merlin 3.0} crate, which implements 
Trevor Perrin's STROBE-based transcript construction. Merlin transcripts 
provide domain-separated, sequential absorption of public inputs and 
prover messages, plus uniform extraction of challenges from the running 
transcript state. Using Merlin uniformly across all six proving systems 
ensures that no two proofs share a challenge by accident and that the 
implementation matches the strong Fiat-Shamir variant assumed by the paper.

\paragraph{Hashing and randomness.}
SHA-256 (via the \texttt{sha2} crate) is used for general-purpose hashing 
outside Merlin transcripts. Randomness flows through the standard 
\texttt{rand\_core} trait, with \texttt{rand\_chacha} as a deterministic 
seedable RNG for tests. Secret material is wrapped with the \texttt{zeroize} 
crate so that secrets are erased from memory when their containers are 
dropped --- a defense-in-depth measure that does not affect protocol 
correctness but reduces the impact of memory-disclosure bugs.

\paragraph{Optional dependencies.}
A set of dependencies is gated behind the \texttt{lib-proofs} Cargo 
feature: \texttt{bulletproofs}, \texttt{sigma-protocols}, \texttt{zkp}, 
\texttt{one-of-many-proofs}, and the older crypto stack (\texttt{curve25519-dalek-ng}, 
\texttt{merlin 2}, \texttt{rand 0.7}) needed for byte-level bridging. 
These are used only by the library-backed proof implementations described 
in Section~\ref{sec:impl-two-paths}, never by the protocol itself.

\section{Repository Structure}
\label{sec:impl-layout}

The repository follows the standard Cargo layout:

\begin{verbatim}
src/
  proofs/                Six proving systems (representation, DLEQ,
                         encryption validity, serial validity, bit
                         vector, commitment set), one file per proof.
    lib_backed/          Optional library-backed equivalents of the
                         six proofs (feature `lib-proofs`).
    bridge.rs            Byte-level bridge between the home crypto
                         stack (dalek 4.1 + merlin 3) and library
                         stacks living on older versions.
  elgamal/               Threshold ElGamal: encryption, partial
                         decryption, and BSGS lookup for recovering
                         the small-integer plaintext from g^T.
  dkg/                   Pedersen distributed key generation with
                         Feldman commitments.
  election/              Protocol orchestration: SetupElection,
                         SetupTally, SetupVoter, Vote, VerifyBallot,
                         Tally, VerifyTally, plus the bulletin board.
  signature.rs           Context-bound Schnorr signatures used to
                         authenticate organizer, tallier, and voter
                         messages.
benches/                 Criterion benchmark suites: per-proof
                         microbenchmarks, full election runs at ten
                         voter scales, proof-size reports, and the
                         hand-rolled vs library comparison harness.
tests/                   Integration tests, including a full
                         end-to-end election and a coercion-resistance
                         re-vote scenario.
vendor/
  one-of-many-proofs/    Vendored fork of cargodog/one-of-many-proofs,
                         used by `lib_backed::commitment_set`. Patched
                         to drop the `simd_backend` and `nightly`
                         features (which depended on packed_simd_2,
                         no longer building on stable rustc).
                         MIT-licensed; original by cargodog.
thesis/                  This document.
\end{verbatim}

\section{Module: Threshold Key Generation}
\label{sec:impl-dkg}

The \texttt{dkg/} module implements Pedersen-style distributed key 
generation with Feldman verifiable secret sharing, exactly as specified 
in Phase 1 of the protocol (Section~\ref{sec:phase-1-dkg}).

The module exposes two types: a \texttt{TallierState} that drives one 
authority through the protocol, and a top-level \texttt{run\_dkg} 
function that orchestrates the full $\eta$-party interaction in 
simulation (production deployments would replace this with a real 
networking layer). Each authority's state machine carries:

\begin{itemize}
    \item Its own secret polynomial $f_\alpha$ of degree $t - 1$, sampled 
    by drawing $t$ random scalars from the configured RNG.
    \item Its Feldman commitment vector 
    $\mathbf{C}_\alpha = (g^{a_{\alpha, 0}}, \ldots, g^{a_{\alpha, t-1}})$, 
    computed once and broadcast.
    \item A NIZK proof of knowledge of $a_{\alpha, 0}$ (a representation 
    proof from \texttt{src/proofs/representation.rs}), which prevents 
    key-cancellation attacks.
    \item Outgoing private shares $f_\alpha(\beta)$ for every other 
    authority $\beta$.
    \item Verified incoming shares from every other authority.
\end{itemize}

Share verification follows the exact equation 
$g^{f_\beta(\alpha)} = \prod_{j=0}^{t-1} (C_{\beta, j})^{\alpha^j}$ from 
the paper, evaluated in constant time using \texttt{curve25519-dalek}'s 
fixed-base multiplication. A failing share aborts the DKG and surfaces 
which authority's contribution was malformed.

The final share $sk_\alpha = \sum_\beta f_\beta(\alpha)$ and the public 
key $pk = \prod_\beta C_{\beta, 0}$ are derived once the receive phase 
completes. Both are returned alongside a NIZK of correctness for 
$pk_\alpha = g^{sk_\alpha}$, which becomes part of the public election 
transcript.

\section{Module: Voter Registration}
\label{sec:impl-registration}

The \texttt{election/setup.rs} file handles voter registration. Per 
Section~\ref{sec:phase-2-registration}, each voter generates a private 
seed $s_i$ and blinding factor $r_i$, both as fresh Ristretto scalars 
drawn from the RNG, and computes her Pedersen commitment 
$\text{cm}_i = g^{s_i} h^{r_i}$.

The voter additionally produces a representation proof 
$\pi_{\text{rep}, i}$ asserting knowledge of an opening $(s_i, r_i)$ 
of $\text{cm}_i$, using the proof in \texttt{src/proofs/representation.rs}. 
The triple $(\text{cm}_i, \pi_{\text{rep}, i}, \text{vk}_i)$ (where $\text{vk}_i$ is the voter's signing-verification key) 
is signed with the voter's long-term signing key (handled in 
\texttt{src/signature.rs}) and posted to the bulletin board.

The implementation does not use a Merkle tree over the voter list. 
This matches the paper: the eligibility proof inside each ballot 
(Section~\ref{sec:impl-commitment-set}) is the Bootle et al.\ 
commitment-set proof, which already produces $O(\log n)$-size proofs 
operating directly on the flat list 
$L = (\text{cm}_1, \ldots, \text{cm}_n)$. A Merkle layer would add 
authentication-path bookkeeping without reducing proof size or 
verification cost. The bulletin-board representation is therefore a 
flat \texttt{Vec<RistrettoPoint>}.

\section{Module: Ballot Construction}
\label{sec:impl-ballot}

Ballot construction lives in \texttt{src/election/ballot.rs} and pulls 
together five of the six proving systems. For each ballot, the voter 
performs the steps of Section~\ref{sec:phase-5-encryption} through 
Section~\ref{sec:phase-7-publication}.

\subsection{Encryption}
\label{sec:impl-encryption}

The \texttt{elgamal/} module provides the threshold ElGamal interface. 
For each candidate slot $j = 1, \ldots, k'$ (after the padding to length 
$k' = k + \ell_{\max} - \ell_{\min}$ described in 
Section~\ref{sec:phase-5-encryption}), the voter:

\begin{enumerate}
    \item Samples a fresh nonce $\rho_{i, j}$.
    \item Computes the ElGamal pair 
    $E_{i, j} = (g^{\rho_{i, j}}, g^{v_{i, j}} \cdot pk^{\rho_{i, j}})$, 
    where multiplication and exponentiation are Ristretto group operations 
    via \texttt{curve25519-dalek}.
    \item Produces an encryption-validity proof from 
    \texttt{src/proofs/enc\_validity.rs}, asserting that the pair encrypts 
    a known scalar under $pk$.
\end{enumerate}

\subsection{Bit-vector validity proof}
\label{sec:impl-bit-vector}

The bit-vector proof at \texttt{src/proofs/bit\_vector.rs} implements the 
generalization of Bootle et al.'s ring-signature bit proof described in 
Appendix A of the Aura paper. It produces a single $\Sigma$-protocol 
proof asserting two facts simultaneously:

\begin{itemize}
    \item Each $v_{i, j} \in \{0, 1\}$.
    \item $\sum_{j=1}^{k'} v_{i, j} = \ell_{\max}$ (a single fixed sum, 
    after padding).
\end{itemize}

The implementation follows Appendix A line-by-line: the prover commits 
to vectors $A$, $C$, $D$ derived from random masks and the bit values, 
the verifier challenge $x$ is extracted from a Merlin transcript, and 
the responses $\{f_i\}, z_A, z_C$ are computed and packed into the 
proof object. Verification evaluates the two algebraic identities of 
Appendix A using a single multi-scalar multiplication (Pippenger), 
keeping per-proof verifier work near-linear in $k'$.

\subsection{Commitment-set (one-of-many) proof}
\label{sec:impl-commitment-set}

The commitment-set proof at \texttt{src/proofs/commitment\_set.rs} is 
the largest and most subtle component of the implementation. It proves 
that there exists an index $\ell$ and randomness $r$ such that 
$\text{cm}_\ell - C'_i = h^r$, where $C'_i = g^{s_i} h^{r'_i}$ is the 
voter's freshly-randomized commitment to her seed.

The construction follows Bootle et al.\ (ESORICS 2015) with the Spark 
modification of Jivanyan and Feickert~\cite{lelantus-spark}: the voter's 
secret index $\ell$ is decomposed into $m$ digits in base $n$, where 
$n^m = N_{\text{voters}}$. The proof commits to each digit, runs $m$ 
parallel Schnorr-style proofs that each digit lies in $\{0, 1, \ldots, n-1\}$, 
and uses a polynomial-evaluation argument to prove that the public 
commitments line up with the secret index.

The parameters $(n, m)$ are configurable: the implementation uses 
$n = 4$ by default (a reasonable trade-off between proof size and 
verifier work), so that $N = 4^m$ for $m = 2, 3, 4, 5, 6, 7, 8, 9, 10$ 
covers election sizes from $16$ to over a million. Empirical proof 
sizes are tabulated in Chapter~\ref{ch:results}.

The Fiat-Shamir transcript ordering is critical: the verifier challenge 
must be extracted from a transcript that has absorbed every public input 
of the proof (the commitments $\{C_i\}$, the offset $C'$, the digit 
commitments) before any prover messages. The implementation uses a 
single Merlin transcript per ballot, threaded through every sub-proof, 
which guarantees that no challenge is ever extracted before its 
preceding messages have been absorbed.

\subsection{Serial validity proof}
\label{sec:impl-serial-validity}

The encrypted serial number $(D'_i, E'_i) = (g^{r''_i}, u^{s_i} \cdot pk^{r''_i})$ 
is produced in \texttt{src/election/ballot.rs} alongside the rest of the 
ballot. Per Section~\ref{sec:phase-4-serial}, this is what enables 
Aura's coercion resistance.

The serial-validity proof at \texttt{src/proofs/serial\_validity.rs} 
asserts that $(D'_i, E'_i)$ encrypts $u^{s_i}$ for the same $s_i$ that 
opens $C'_i$. Mechanically, this is a simultaneous representation 
proof in three independent generator triples, compiled through Merlin 
into a single non-interactive proof. Without this binding, a malicious 
voter could register one seed but encrypt a different value as her 
serial number, defeating coercion resistance.

\subsection{Ballot assembly and signing}
\label{sec:impl-ballot-assembly}

The complete ballot is the tuple
\[
    B_i = \big(\{(D_{i, j}, E_{i, j}, \pi_{\text{enc}, i, j})\}_{j=1}^{k'},\; 
    \pi_{\text{bit}, i},\; C'_i,\; D'_i,\; E'_i,\; \pi_{\text{set}, i}\big),
\]
plus a serial-validity proof $\pi_{\text{ser}, i}$ that binds $B_i$ as 
its transcript. The voter signs $B_i \| \pi_{\text{ser}, i}$ with her 
long-term signing key (using the context-bound Schnorr signature scheme 
from \texttt{src/signature.rs}), and submits the signed package to the 
bulletin board.

\section{Module: Verification}
\label{sec:impl-verification}

\texttt{src/election/verify.rs} implements the verifier side of 
Section~\ref{sec:phase-7-publication}. For each ballot on the bulletin 
board, the verifier:

\begin{enumerate}
    \item Checks the voter signature on $B_i$ using the voter's 
    verification key from $L_{\text{voters}}$.
    \item Checks that the ballot is not a byte-level duplicate of an 
    already-accepted ballot.
    \item Verifies the serial-validity proof $\pi_{\text{ser}, i}$ 
    against the public inputs of the ballot, using $B_i$ itself as the 
    transcript binding.
    \item Verifies each per-slot encryption-validity proof 
    $\pi_{\text{enc}, i, j}$.
    \item Verifies the bit-vector proof $\pi_{\text{bit}, i}$ on the 
    aggregated ciphertext $\sum_j E_{i, j}$.
    \item Verifies the commitment-set proof $\pi_{\text{set}, i}$ 
    against the flat voter list $L$ and the offset $C'_i$.
\end{enumerate}

A ballot that fails any check is rejected and not included in the tally. 
Every check is implemented as variable-time, since the verifier only 
operates on public data; constant-time would be a needless cost.

The verifier supports \emph{batch verification} of compatible proofs. 
Encryption-validity proofs across all ballots can be combined into a 
single multi-scalar multiplication using random linear combination, 
amortizing per-proof verifier work as the batch grows. The same trick 
applies to the representation proofs in DKG and to the DLEQ proofs 
produced during decryption. Cross-ballot batch verification is 
benchmarked in \texttt{benches/election\_n*.rs} as 
\texttt{verify\_batch\_encval} groups, comparing individual versus 
batched verification at batch sizes 1, 10, 50, and 100.

\section{Module: Tally and Threshold Decryption}
\label{sec:impl-tally}

The tally module \texttt{src/election/tally.rs} implements the two-pass 
tally of Section~\ref{sec:phase-9-decryption}.

\paragraph{Pass 1: serial decryption and deduplication.}
For each accepted ballot, a threshold cohort of $t$ talliers each 
computes a partial decryption $R_{\alpha, i} = (D'_i)^{sk_\alpha}$ of 
the ballot's encrypted serial number, using the partial-decryption 
function in \texttt{src/elgamal/decrypt.rs}. Each tallier additionally 
produces a DLEQ proof (\texttt{src/proofs/dleq.rs}) asserting that the 
same $sk_\alpha$ that produced her published $pk_\alpha$ was used to 
compute $R_{\alpha, i}$. The combiner verifies all DLEQ proofs, then 
reconstructs $(D'_i)^{sk}$ via Lagrange interpolation:
\[
    (D'_i)^{sk} = \prod_{\alpha \in S} R_{\alpha, i}^{\lambda_\alpha},
\]
and recovers $u^{s_i} = E'_i \cdot \big((D'_i)^{sk}\big)^{-1}$. Ballots 
sharing the same $u^{s_i}$ are grouped, and all but the most recent 
(by bulletin board position) are discarded. This realizes coercion 
resistance: a voter who recast her ballot during the election has only 
the most recent ballot counted.

\paragraph{Pass 2: aggregate decryption.}
For each candidate $j$, the homomorphic aggregate ciphertext 
$E_\Sigma^{(j)} = \prod_i E_{i, j}$ is computed over the surviving 
ballots --- a pure public-arithmetic step that uses no key material. 
The same threshold cohort then partial-decrypts each of the $k$ 
aggregates, and Lagrange interpolation produces the group element 
$g^{T_j}$ for each candidate.

\paragraph{BSGS lookup.}
Recovering $T_j$ from $g^{T_j}$ requires solving a discrete log in a 
known small range. The implementation uses a precomputed 
\emph{baby-step giant-step} table over the range $[0, n]$, located in 
\texttt{src/elgamal/decrypt.rs::bsgs}. For an election with $n$ voters, 
the table has $\sqrt{n}$ entries and the lookup costs $O(\sqrt{n})$ time, 
which is negligible compared to the cost of producing the aggregate 
ciphertext in the first place. The recovered tuple 
$(T_1, \ldots, T_k)$ is the official election result.

\section{Two Implementations of the Proofs}
\label{sec:impl-two-paths}

For each of the six proving systems, the repository contains two 
parallel implementations:

\begin{enumerate}
    \item \textbf{Hand-rolled} (\texttt{src/proofs/}, default). Reproduces 
    the Aura paper's constructions directly: same statements, same 
    transcript layout, same generators. This is what the protocol uses 
    for every election, ballot, DKG, and tally operation, regardless of 
    feature flags.
    \item \textbf{Library-backed} (\texttt{src/proofs/lib\_backed/}, 
    feature \texttt{lib-proofs}). Each proof delegates to an off-the-shelf 
    Rust ZK library running on its own crypto stack, with byte-level 
    bridging at the module boundary.
\end{enumerate}

The mapping between proofs and libraries is shown in 
Table~\ref{tab:impl-lib-mapping}.

\begin{table}[H]
\centering
\caption{Library mapping for the alternative implementations of each 
proving system. ``Bridged'' means the library lives on a different 
curve/transcript stack than Aura's home (\texttt{curve25519-dalek 4.1} 
+ \texttt{merlin 3}), with a 32-byte canonical encoding bridge at the 
module boundary.}
\label{tab:impl-lib-mapping}
\small
\begin{tabular}{p{0.4cm}p{3.0cm}p{3.5cm}p{4.5cm}p{1cm}}
\toprule
\# & Proof & Library & Crypto stack & Bridged \\
\midrule
1 & Representation & \texttt{zkp 0.8} & dalek-ng 3 / merlin 2 / rand 0.7 & yes \\
2 & DLEQ & \texttt{sigma-protocols 0.5} & dalek 4.1 & no \\
3 & EncryptionValidity & \texttt{zkp 0.8} & dalek-ng 3 / merlin 2 / rand 0.7 & yes \\
4 & SerialValidity & \texttt{zkp 0.8} & dalek-ng 3 / merlin 2 / rand 0.7 & yes \\
5 & BitVector & \texttt{bulletproofs} (R1CS) & dalek 4.1 / merlin 3 & no \\
6 & One-of-N & \texttt{one-of-many-proofs} (vendored) & dalek 2 / merlin 2 / rand 0.7 & yes \\
\bottomrule
\end{tabular}
\end{table}

The motivation for maintaining both is to surface what a developer 
would actually pay --- in performance, proof size, and protocol fidelity 
--- if they tried to assemble Aura on top of widely-used Rust ZK libraries 
instead of implementing the $\Sigma$-protocols by hand. Two of the six 
library-backed proofs do not prove the exact statement of the 
hand-rolled paper construction:

\begin{itemize}
    \item \textbf{Bit-vector (bulletproofs R1CS).} Certifies the predicate 
    ($b_i \in \{0, 1\}$, $\sum_i b_i = w$) using bulletproofs' default 
    \texttt{PedersenGens}, but without binding to Aura's vector commitment 
    $B = r \cdot H + \sum_i b_i \cdot G_i$. The hand-rolled proof 
    additionally binds to $B$.
    \item \textbf{One-of-N (\texttt{one-of-many-proofs}).} Uses the 
    library's hardcoded basepoint $G = \texttt{RISTRETTO\_BASEPOINT\_POINT}$ 
    rather than Aura's \texttt{params.h}, only supports set sizes that 
    are exact powers of two, and recomputes $C'$ internally using the 
    library's basepoint. The hand-rolled proof uses Aura's generators 
    and the Bootle/Spark construction.
\end{itemize}

These statement-fidelity caveats are documented in the module docstrings 
and disclosed alongside the timing comparisons in Chapter~\ref{ch:results}.

The library-backed paths are never exercised by the protocol itself: 
all call sites in \texttt{src/election/}, \texttt{src/elgamal/}, and 
\texttt{src/dkg/} import directly from the hand-rolled \texttt{crate::proofs::*} 
modules. The library-backed code is compiled only when the 
\texttt{lib-proofs} feature is enabled, and is exercised exclusively 
by the comparison benchmark suite \texttt{proof\_primitives\_lib} and 
its own unit tests.

\section{Testing}
\label{sec:impl-testing}

The test suite consists of 57 unit tests covering the hand-rolled 
proving systems, ElGamal operations, DKG, and signatures, plus 5 
integration tests in \texttt{tests/}. With the \texttt{lib-proofs} 
feature enabled, an additional 22 unit tests cover the library-backed 
proofs --- one module per proof, one happy path plus several edge cases 
each.

\paragraph{Unit tests.}
Each proving system has tests for: correctness on randomly generated 
witnesses, soundness against tampered proofs (every component bit-flipped 
in turn, expected to fail verification), and serialization round-trips 
where serde support is enabled. The DKG module additionally tests 
abort behavior for malformed Feldman commitments and for shares that 
fail Feldman verification.

\paragraph{Integration tests.}
The \texttt{tests/full\_election.rs} file contains two end-to-end tests:

\begin{itemize}
    \item \texttt{full\_election\_8\_voters\_3\_talliers\_4\_candidates}: 
    runs the entire protocol pipeline on small parameters --- 
    \textsf{SetupElection}, \textsf{SetupTally} (DKG with $\eta = 3$, 
    $t = 2$), \textsf{SetupVoter} for each of 8 voters, ballot 
    construction and verification, two-pass tally. The expected result 
    is asserted byte-for-byte against the actual decrypted tally vector.
    \item \texttt{coercion\_resistance\_ballot\_recast}: a voter casts 
    a ballot, then casts a different ballot for the same election. The 
    test asserts that both ballots are accepted to the bulletin board 
    (since the encrypted serial numbers are indistinguishable during 
    voting), and that after deduplication only the second ballot 
    contributes to the tally.
\end{itemize}

\paragraph{Worked example as regression test.}
The toy parameters of Chapter~\ref{ch:byhand} ($p = 47$, $q = 23$ over 
a small Schnorr group) are not directly executable by the protocol 
implementation, since the Ristretto255 group used by the protocol does 
not admit such tiny instantiations. However, the worked example serves 
as a structural regression test: each protocol phase from 
Chapter~\ref{ch:byhand} maps to a specific module function, and the 
end-to-end integration tests above exercise the same control flow on 
real-sized parameters. Running 
\texttt{cargo test full\_election\_8\_voters} reproduces, on Ristretto255, 
the exact sequence of phases shown by hand on $p = 47$ in 
Chapter~\ref{ch:byhand}.

\paragraph{Continuous validation.}
The \texttt{cargo clippy -- -D warnings} target is run as part of the 
test workflow, ensuring that no Clippy lint regressions are merged. All 
tests are deterministic when run with a fixed seed, making CI failures 
reproducible without timing-related flakiness.
% =============================================================
% CHAPTER 6 — RESULTS
% =============================================================
\chapter{Benchmarks and Results}
\label{ch:results}

This chapter reports performance measurements of the implementation 
described in Chapter~\ref{ch:implementation}. All numbers come from 
deterministic benchmarks executed on the development machine, with the 
exact hardware, software stack, and randomness seed recorded in the 
repository for full reproducibility.

\section{Benchmarking Methodology}
\label{sec:bench-methodology}

\paragraph{Hardware.}
All measurements in the body of this chapter were taken on a single AWS
EC2 instance backed by an Intel Xeon Platinum 8488C (Sapphire Rapids),
exposed to the guest as $2$ vCPU / $2$ threads, with $3.7$ GiB of
system memory. The instance is a KVM/full-virtualisation guest; the L1d
cache is $48$ KiB, L1i $32$ KiB, L2 $2$ MiB, and the L3 (shared with the
host socket) measures $105$ MiB. Benchmarks run single-threaded and use
\texttt{rdtsc}-based wall-clock timing via Criterion. The host operating
system is Ubuntu $24.04.4$ LTS on Linux kernel $6.17.0$ (AWS). A separate
library-backed comparison was performed on a different machine (laptop)
and is reported in Appendix~\ref{app:lib-comparison}; all other benchmarks
in this chapter use the AWS environment described above.

\paragraph{Software stack.}
The implementation was compiled with \texttt{rustc 1.95.0} (released
April~2026) and \texttt{cargo 1.95.0}, in release mode with default
LTO and codegen settings. The cryptographic group is Ristretto255 via
\texttt{curve25519-dalek 4.1}; transcripts use \texttt{merlin 3.0};
hashing uses \texttt{sha2 0.10}. Multi-scalar multiplication uses
\texttt{curve25519-dalek}'s built-in Pippenger implementation in
variable-time mode (the verifier's working data is public).

\paragraph{Determinism.}
Every benchmark seeds its randomness from a fixed \texttt{ChaCha20Rng}
state (seed \texttt{20260412}). Re-running on the same machine produces
means within Criterion's reported confidence intervals. The exact
hardware, software, and Git commit used to produce the numbers in this
chapter are recorded verbatim in \texttt{thesis/results/env.txt} of the
repository (and \texttt{env\_laptop.txt} for the appendix run).

\paragraph{Measurement methodology.}
All benchmarks use Criterion~0.5 with default settings: a warm-up phase 
followed by sample collection until the configured sample count is 
reached. Sample counts are 100--200 for fast microbenchmarks 
(microsecond-scale operations like Schnorr proofs) and 10--50 for slower 
benchmarks (millisecond-scale operations like full ballot construction 
and full election lifecycles). Reported means are accompanied by 
half-confidence-interval percentages drawn from Criterion's internal 
bootstrap analysis. We report the mean rather than median because 
asymmetry between the two never exceeded $5\%$ in any benchmark.

\paragraph{Tier.}
The benchmarks reported in this chapter cover election sizes
$N \in \{16, 64, 256, 1{,}024, 4{,}096, 16{,}384\}$. A further attempt
at $N = 65{,}536$ exhausted the $3.7$ GiB AWS instance memory during
ballot construction and did not produce timing data; we treat that
configuration as a soft ceiling for the hardware used here. Proof-size
measurements (which are independent of runtime cost) cover the full
range up to $N = 10^6$ and are reported in the size tables below.

\paragraph{What is measured.}
The benchmark suite is organized into five logical groups:
\begin{itemize}
\item \textbf{G1} --- Per-proof microbenchmarks for the six hand-rolled 
proving systems, with prove and verify timed separately at multiple 
parameter values.
\item \textbf{G2} --- Byte sizes of every proof and full ballot, measured 
across all supported election scales.
\item \textbf{G3} --- Batch-verification amortization for proofs that 
support it (representation, DLEQ, encryption-validity), at batch sizes 
1, 10, 25, 50, and 100.
\item \textbf{G4} --- Full election lifecycle timings (DKG, voter 
registration, ballot create/verify, tally) at each tier scale.
\item \textbf{G6} --- Tally pipeline breakdown (per-tallier serial 
decryption versus full pipeline including aggregate decryption).
\end{itemize}
A sixth group (G5, hand-rolled vs library-backed comparison) was
measured on a separate laptop machine because the \texttt{lib-proofs}
Cargo feature pulls additional dependencies that were not part of the
AWS run; we report the results separately in
Appendix~\ref{app:lib-comparison}.

\section{Headline numbers}
\label{sec:bench-headline}

Before drilling into per-phase microbenchmarks, the table below
collapses the AWS run into a one-row-per-deployment-scale executive
summary. Three scenarios are picked from the measured range to anchor
intuition: \emph{boardroom} ($N = 256$, e.g.\ a faculty senate vote),
\emph{small organisation} ($N = 4{,}096$, e.g.\ a small university
elections), and \emph{mid-scale} ($N = 16{,}384$, e.g.\ a city-council
ward). Per-phase columns are single-instance wall-clock times taken
verbatim from Table~\ref{tab:bench-lifecycle}; the right-most column
is the end-to-end serialized total assuming a single voter at a time
and a single verifier ($\text{DKG} + N \times \text{registration} +
N \times \text{ballot create} + N \times \text{ballot verify} +
\text{tally}$). Real deployments will parallelise both ballot
construction (one voter per voter device) and verification (a
verifier pool), but the serialized figure is the worst-case ceiling
and the right number to anchor sizing decisions.

\begin{table}[H]
\centering
\caption{Headline performance numbers for three target deployment
scales (AWS run, single-threaded). All per-phase numbers come from
Table~\ref{tab:bench-lifecycle}; the rightmost column is the
serialized end-to-end total. Tally is the full pipeline including
BSGS lookup.}
\label{tab:bench-headline}
\small
\begin{tabular}{lrrrrrr}
\toprule
Scenario & DKG & Voter reg. & Ballot create & Ballot verify & Tally & Serial e2e \\
\midrule
Boardroom ($N = 256$)        & $1.06$ ms & $116$ µs & $13.2$ ms & $4.27$ ms & $200$ ms  & $\approx 4.7$ s   \\
Small org ($N = 4{,}096$)    & $1.06$ ms & $118$ µs & $157$ ms  & $37.4$ ms & $4.33$ s  & $\approx 13.3$ min \\
Mid-scale ($N = 16{,}384$)   & $1.09$ ms & $118$ µs & $668$ ms  & $145$ ms  & $31.8$ s  & $\approx 3.7$ h    \\
\bottomrule
\end{tabular}
\end{table}

The serialized end-to-end figure is dominated by per-voter ballot
construction at every scale --- $\approx 72\%$ of the boardroom
total, $\approx 80\%$ at small-org scale, and $\approx 82\%$ at
mid-scale --- which is the cost paid in parallel by the voter pool,
not by any centralised infrastructure. The non-parallelisable
components (DKG once, full-pipeline tally once) stay under $32$
seconds even at the largest measured scale.

\section{Per-Phase Microbenchmarks}
\label{sec:bench-microbenchmarks}

This section reports the cost of each individual cryptographic operation
in isolation, before they are composed into ballots or full elections.

\subsection{Key Generation Costs}
\label{sec:bench-dkg}

The distributed key generation (Phase 1 of Chapter~\ref{ch:aura}) involves 
all $\eta$ authorities running Pedersen-Feldman VSS interactively. For 
the configuration tested ($\eta = 3$, $t = 2$), DKG is essentially 
constant in election size $N$ --- it depends only on the number of 
authorities and the threshold, not on the voter count. Measured
wall-clock times across the AWS run are:
$1.130$ ms at $N = 16$, $1.065$ ms at $N = 64$, $1.060$ ms at $N = 256$,
$1.062$ ms at $N = 1{,}024$, $1.060$ ms at $N = 4{,}096$, and
$1.091$ ms at $N = 16{,}384$ (mean $\pm 0.5\%$ across runs). The
variation across $N$ is well within measurement noise; DKG is
genuinely $O(\eta^2 \cdot t)$ in the authority parameters, independent
of voters. Figure~\ref{fig:dkg-vs-n} confirms the constancy visually:
the dashed reference line is the mean over all six measurements, and
no individual measurement deviates by more than $\sim 7\%$.

\begin{figure}[H]
\centering
\includegraphics[width=0.75\textwidth]{dkg_vs_n.pdf}
\caption{Distributed-key-generation cost is independent of voter set
size $N$ (AWS run, $\eta = 3$ authorities, threshold $t = 2$).
Error bars are Criterion $95\%$ confidence intervals.
\figdata{thesis/results/criterion/N\{X\}\_dkg/2of3/new/estimates.json}}
\label{fig:dkg-vs-n}
\end{figure}

This is consistent with the protocol design: each authority performs
$O(t)$ scalar multiplications for its own polynomial commitments, $O(\eta)$ 
share evaluations, and $O(\eta \cdot t)$ verifications of received shares. 
For the small-cohort case ($\eta \le 10$) typical of real elections, 
DKG is a one-time cost in the low milliseconds.

\subsection{Schnorr-Type Proof Costs}
\label{sec:bench-schnorr}

The four Schnorr-style proving systems (representation, DLEQ, 
encryption-validity, serial-validity) take constant time per proof, with 
no dependency on election parameters. Their measured prove and verify 
times are summarized in Table~\ref{tab:bench-schnorr}.

\begin{table}[H]
\centering
\caption{Per-proof prove and verify times for the four Schnorr-style 
proving systems. Times are mean $\pm$ half-confidence-interval percentage 
(Criterion).}
\label{tab:bench-schnorr}
\small
\begin{tabular}{lcc}
\toprule
Proof system & Prove & Verify \\
\midrule
Representation ($n=1$ generator)  & $43.70$ µs $\pm 0.5\%$ & $74.73$ µs $\pm 0.5\%$ \\
Representation ($n=2$ generators) & $57.53$ µs $\pm 0.4\%$ & $83.37$ µs $\pm 0.5\%$ \\
Representation ($n=4$ generators) & $86.85$ µs $\pm 0.5\%$ & $101.7$ µs $\pm 0.4\%$ \\
Representation ($n=8$ generators) & $139.4$ µs $\pm 0.4\%$ & $138.0$ µs $\pm 0.5\%$ \\
DLEQ                              & $85.71$ µs $\pm 0.4\%$ & $146.3$ µs $\pm 0.5\%$ \\
Encryption validity               & $120.4$ µs $\pm 0.5\%$ & $177.9$ µs $\pm 0.4\%$ \\
Serial validity                   & $191.5$ µs $\pm 0.4\%$ & $280.6$ µs $\pm 0.4\%$ \\
\bottomrule
\end{tabular}
\end{table}

The four times scale with the number of group operations in the underlying
relation: representation is the simplest (two generators, two scalars),
DLEQ adds a second relation, encryption-validity proves an ElGamal
encryption (three generators), and serial-validity is the most complex
(four group elements, three witnesses, two relations). All four operations
sit comfortably in the sub-millisecond range; even on commodity hardware,
a voter constructing a ballot with these proofs adds well under a
millisecond of total Schnorr work.

\begin{figure}[H]
\centering
\includegraphics[width=0.85\textwidth]{g1_proof_primitives.pdf}
\caption{Per-proof prover vs. verifier cost at the largest scale
measured for each proof family (AWS run), plotted on a log y-axis.
Representation ($n = 8$), DLEQ, encryption-validity and
serial-validity all stay in the low hundreds of µs; bit-vector
($k' = 64$) and commitment-set ($N = 256$) are an order of magnitude
heavier. \figdata{thesis/results/timings.csv (source=aws)}}
\label{fig:g1-proof-primitives}
\end{figure}

\begin{figure}[H]
\centering
\includegraphics[width=0.95\textwidth]{g1_proof_primitives_scaling.pdf}
\caption{Log--log scaling of the three parameterised proof primitives.
\textbf{Left:} representation proof prove/verify times across $n \in \{1,2,4,8\}$
generators. \textbf{Middle:} bit-vector proof across padded slot counts
$k' \in \{4,5,8,12,16,24,32,48,64\}$. \textbf{Right:} commitment-set
proof across $N \in \{4, 16, 64, 256, 1{,}024, 4{,}096\}$. All three
families show clean linear or near-linear scaling on log axes.
\figdata{thesis/results/timings.csv (source=aws)}}
\label{fig:g1-scaling}
\end{figure}

\begin{figure}[H]
\centering
\includegraphics[width=0.85\textwidth]{representation_violin.pdf}
\caption{Criterion violin plot for the representation proof, showing the
distribution of measured samples at each generator count $n$
(\texttt{prove}). The tight density at every $n$ confirms low
measurement noise on the AWS instance.
\figdata{thesis/results/criterion/representation/report/violin.svg}}
\label{fig:crit-representation}
\end{figure}

\subsection{Bit-Vector Proof Costs}
\label{sec:bench-bitvec}

The bit-vector proof (Section~\ref{sec:phase-6-validity}, statement P3) 
asserts that an encrypted bit vector contains only zeros and ones, with 
a fixed weight. Its cost scales linearly with the padded slot count $k'$.

\begin{table}[H]
\centering
\caption{Bit-vector proof prove and verify times as a function of padded 
slot count $k'$.}
\label{tab:bench-bitvec}
\small
\begin{tabular}{ccc}
\toprule
$k'$ & Prove & Verify \\
\midrule
4   & $243.3$ µs $\pm 0.4\%$ & $193.7$ µs $\pm 0.6\%$ \\
5   & $274.5$ µs $\pm 0.5\%$ & $206.3$ µs $\pm 0.5\%$ \\
8   & $375.0$ µs $\pm 0.5\%$ & $248.6$ µs $\pm 0.5\%$ \\
12  & $503.8$ µs $\pm 0.4\%$ & $306.4$ µs $\pm 0.5\%$ \\
16  & $636.8$ µs $\pm 0.5\%$ & $362.6$ µs $\pm 0.5\%$ \\
24  & $895.9$ µs $\pm 0.4\%$ & $480.4$ µs $\pm 0.4\%$ \\
32  & $1.163$ ms $\pm 0.4\%$ & $613.5$ µs $\pm 0.5\%$ \\
48  & $1.730$ ms $\pm 2.3\%$ & $876.6$ µs $\pm 0.4\%$ \\
64  & $2.203$ ms $\pm 0.5\%$ & $1.131$ ms $\pm 0.4\%$ \\
\bottomrule
\end{tabular}
\end{table}

The doubling pattern is visible: from $k' = 4$ to $k' = 64$ (a $16\times$
increase), prove time grows from $243$ µs to $2.20$ ms ($9.1\times$),
and verify time grows from $194$ µs to $1.13$ ms ($5.8\times$). The
sub-linear growth on the verifier side reflects the multi-scalar
multiplication used at the verification step, which has more aggressive
amortization than a naive sum.

\begin{figure}[H]
\centering
\includegraphics[width=0.85\textwidth]{bit_vector_violin.pdf}
\caption{Criterion violin plot for the bit-vector proof across padded
slot counts $k' \in \{4, 5, 8, 12, 16, 24, 32, 48, 64\}$ (\texttt{prove}).
The roughly linear envelope is consistent with the $O(k')$ scaling
reported in Table~\ref{tab:bench-bitvec}.
\figdata{thesis/results/criterion/bit\_vector/report/violin.svg}}
\label{fig:crit-bit_vector}
\end{figure}

\subsection{Commitment-Set (Eligibility) Proof Costs}
\label{sec:bench-commset}

The commitment-set proof (Section~\ref{sec:phase-6-validity}, statement 
P1) is the largest individual component of an Aura ballot. It is also 
the most parameter-sensitive: its cost depends on both the voter list 
size $N$ and the $n$-ary digit base used in the Bootle-Spark construction.

\begin{table}[H]
\centering
\caption{Commitment-set proof prove and verify times for binary
($n = 2$) versus quaternary ($n = 4$) digit base, across the full
range of $N$ at which the standalone bench was run. The quaternary
variant approximately halves prove time at essentially the same
verify time. The commitment-set proof at $N = 16{,}384$ is captured
indirectly via the ballot-creation cost in
Table~\ref{tab:bench-lifecycle} (where it is the dominant component);
no standalone $N = 16{,}384$ commitment-set bench was run.}
\label{tab:bench-commset}
\small
\begin{tabular}{cccc}
\toprule
$N$ & Variant & Prove & Verify \\
\midrule
4      & $n = 2$ & $463.2$ µs $\pm 0.3\%$ & $95.84$ µs $\pm 0.3\%$ \\
16     & $n = 2$ & $1.164$ ms $\pm 0.6\%$ & $230.8$ µs $\pm 0.3\%$ \\
64     & $n = 2$ & $3.251$ ms $\pm 0.3\%$ & $740.0$ µs $\pm 0.3\%$ \\
64     & $n = 4$ & $2.115$ ms $\pm 0.4\%$ & $694.2$ µs $\pm 0.3\%$ \\
256    & $n = 2$ & $13.39$ ms $\pm 0.4\%$ & $2.677$ ms $\pm 0.3\%$ \\
256    & $n = 4$ & $7.123$ ms $\pm 0.3\%$ & $2.541$ ms $\pm 0.2\%$ \\
1{,}024  & $n = 2$ & $47.16$ ms $\pm 0.3\%$ & $9.276$ ms $\pm 0.2\%$ \\
1{,}024  & $n = 4$ & $22.94$ ms $\pm 0.2\%$ & $8.886$ ms $\pm 1.2\%$ \\
4{,}096  & $n = 2$ & $173.1$ ms $\pm 0.2\%$ & $35.60$ ms $\pm 0.2\%$ \\
\bottomrule
\end{tabular}
\end{table}

Two observations dominate this table.

First, the prove and verify times grow cleanly with $N$, but at very
different rates. Going from $N = 4$ to $N = 4{,}096$ (a $1024\times$
increase in voter count), prove time grows from $0.46$ ms to $173$ ms
(a $\approx 374\times$ increase), and verify time grows from $0.10$ ms
to $35.6$ ms ($\approx 372\times$). Both are slower than the
asymptotic claim of $O(\log N)$
in proof size --- which is only about the byte size of the proof, not 
the cost of producing it. The actual prover and verifier work in the 
Bootle construction grows as $O(N)$ in the worst case, since both 
parties must touch every commitment in the list at least once. What 
the GK construction reduces is the \emph{communication}, not the 
overall computation. This is consistent with what is reported in the 
literature.

\begin{figure}[H]
\centering
\includegraphics[width=0.85\textwidth]{g4_setproof_vs_n.pdf}
\caption{Commitment-set proof prove and verify times as a function of
voter set size $N$, on a log-log scale (AWS run, binary digit base
$n = 2$). The near-linear lines indicate close-to-$O(N)$ growth, with
verification cheaper than proving across the whole range.
\figdata{thesis/results/timings.csv (group=commitment\_set)}}
\label{fig:g4-setproof-vs-n}
\end{figure}

\begin{figure}[H]
\centering
\includegraphics[width=0.85\textwidth]{commitment_set_violin.pdf}
\caption{Criterion violin plot for the commitment-set proof prover
across all measured $N$ on the AWS run. The per-sample dispersion is
tight; the medians span more than two orders of magnitude as $N$
grows from $4$ to $4{,}096$.
\figdata{thesis/results/criterion/commitment\_set/report/violin.svg}}
\label{fig:crit-commitment_set}
\end{figure}

Second, the choice of $n$-ary digit base meaningfully affects prover
cost. Compared at the same $N$, the $n = 4$ variant approximately halves
prove time: $7.12$ ms versus $13.39$ ms at $N = 256$, $22.94$ ms versus
$47.16$ ms at $N = 1{,}024$. Verify times are essentially identical. The
intuition is that with $n = 4$ the prover commits to $\log_4 N$ 
quaternary digits instead of $\log_2 N$ binary digits, halving the 
commitment-and-bit-proof rounds even though each round becomes slightly 
more expensive. Verification work is dominated by the same multi-scalar 
multiplication regardless of base, so it is largely unchanged.

\subsection{Verification Costs}
\label{sec:bench-verify}

A full ballot verification combines the per-proof verifications above. 
The dominant component is the commitment-set proof verification; the 
encryption-validity, serial-validity, and bit-vector verifications add 
a few hundred microseconds each. Full ballot verification times are 
reported in Section~\ref{sec:bench-lifecycle} alongside the rest of the 
election lifecycle.

A separate question is whether \emph{batch} verification of multiple 
ballots improves over verifying them individually. The Aura paper claims 
that batch verification of compatible Schnorr-type proofs amortizes per-proof 
cost into a single multi-scalar multiplication. Our measurements 
(Table~\ref{tab:bench-batch}) confirm this for all three batchable 
proof systems (representation, DLEQ, encryption-validity).

\begin{table}[H]
\centering
\caption{Batch-verification amortization. ``Individual'' is the time to 
verify each proof one at a time; ``Batched'' is the time for a single 
combined multi-scalar multiplication over all proofs at once. Speedup 
is the ratio.}
\label{tab:bench-batch}
\small
\begin{tabular}{cccccc}
\toprule
Proof system & Batch size & Individual & Batched & Speedup \\
\midrule
Representation & 1   & $80.83$ µs & $58.06$ µs & $1.39\times$ \\
Representation & 10  & $797.2$ µs & $430.7$ µs & $1.85\times$ \\
Representation & 25  & $2.027$ ms & $1.044$ ms & $1.94\times$ \\
Representation & 50  & $4.041$ ms & $2.012$ ms & $2.01\times$ \\
Representation & 100 & $8.166$ ms & $3.700$ ms & $2.21\times$ \\
\midrule
DLEQ           & 1   & $145.2$ µs & $80.54$ µs & $1.80\times$ \\
DLEQ           & 10  & $1.442$ ms & $631.2$ µs & $2.28\times$ \\
DLEQ           & 25  & $3.585$ ms & $1.550$ ms & $2.31\times$ \\
DLEQ           & 50  & $7.188$ ms & $2.812$ ms & $2.56\times$ \\
DLEQ           & 100 & $14.38$ ms & $5.249$ ms & $2.74\times$ \\
\midrule
Enc-validity   & 1   & $178.8$ µs & $94.22$ µs & $1.90\times$ \\
Enc-validity   & 10  & $1.772$ ms & $757.3$ µs & $2.34\times$ \\
Enc-validity   & 25  & $4.446$ ms & $1.840$ ms & $2.42\times$ \\
Enc-validity   & 50  & $8.936$ ms & $3.268$ ms & $2.73\times$ \\
Enc-validity   & 100 & $17.92$ ms & $6.029$ ms & $2.97\times$ \\
\bottomrule
\end{tabular}
\end{table}

The speedup grows with batch size, and reaches $2$--$3\times$ at batch
size $100$ across all three proof systems. The intuition is that random
linear combination compresses many independent verification equations
into a single multi-scalar multiplication, after which Pippenger
amortization makes the marginal cost of each additional proof very
small. For a real election with thousands of ballots, batch verification
is the difference between minutes and tens of seconds of total verifier
work.

\begin{figure}[H]
\centering
\includegraphics[width=0.85\textwidth]{g3_batch_amortization.pdf}
\caption{Batch-verification speedup (individual ÷ batched) as a
function of batch size on the AWS run. All three batchable proof
systems show monotonically increasing speedup, plateauing in the
$2$--$3\times$ range at batch size $100$.
\figdata{thesis/results/timings.csv (group=batch\_verify\_*)}}
\label{fig:g3-batch-amortization}
\end{figure}

\begin{figure}[H]
\centering
\includegraphics[width=0.85\textwidth]{batch_verify_encval_lines.pdf}
\caption{Criterion's per-batch ``lines'' plot for the encryption-validity
batch verifier, showing measured time as a function of batch size with
\texttt{individual} (orange) and \texttt{batch} (blue) lines on the same
axes.
\figdata{thesis/results/criterion/batch\_verify\_encval/report/lines.svg}}
\label{fig:crit-batch_verify_encval}
\end{figure}

\paragraph{Batch verification at election scale.}
The amortization above is measured on standalone proofs with no
election context. A more deployment-relevant question is: how long
does it take to batch-verify $k$ complete ballots, each one referring
to an $N$-voter commitment list? Table~\ref{tab:bench-election-batch}
reports that measurement for $N \in \{16, 64, 256, 1{,}024, 4{,}096,
16{,}384\}$ and batch sizes $k \in \{1, 10, 50, 100\}$. The
per-ballot cost in this table is dominated by the commitment-set
verifier inside each ballot, so the total time grows roughly linearly
in both $N$ and $k$ --- batch verification compresses the
\emph{constant-overhead} portion of each verification, but the
$O(N)$ commitment-set work cannot be folded into a single
multi-scalar multiplication.

\begin{table}[H]
\centering
\caption{Election-scale full-ballot batch verification (AWS run):
wall-clock time to verify $k$ complete ballots at once for a voter
list of size $N$. Each row is one election scale; each column is one
batch size; empty cells mark configurations that were not executed.}
\label{tab:bench-election-batch}
\small
\begin{tabular}{lrrrr}
\toprule
$N$ & $k = 1$ & $k = 10$ & $k = 50$ & $k = 100$ \\
\midrule
16       & $1.623$ ms & $16.43$ ms & ---        & ---        \\
64       & $2.155$ ms & $21.77$ ms & $109.5$ ms & ---        \\
256      & $4.321$ ms & $43.28$ ms & $216.1$ ms & $439.4$ ms \\
1{,}024  & $10.62$ ms & $107.3$ ms & $537.9$ ms & $1.073$ s  \\
4{,}096  & $37.14$ ms & $374.4$ ms & $1.861$ s  & $3.732$ s  \\
16{,}384 & $145.4$ ms & $1.456$ s  & $7.301$ s  & $14.58$ s  \\
\bottomrule
\end{tabular}
\end{table}

Two practical points fall out of this table. First, the
\emph{column-wise} ratio is close to constant: doubling $N$ roughly
doubles per-row time across every batch size, confirming that the
commitment-set verifier inside each ballot is what scales with $N$.
Second, the \emph{row-wise} per-ballot cost (each cell divided by its
$k$) drops modestly as $k$ grows --- at $N = 16{,}384$, single-ballot
verification is $145$ ms but per-ballot cost at $k = 100$ is
$\approx 146$ ms, essentially unchanged. The Schnorr-style batch
amortization is real but masked at large $N$ because the
commitment-set $O(N)$ work cannot be batched, only the constant
per-ballot Schnorr work can. This is the central practical reason a
$10^4$-voter deployment will deploy a parallel verifier pool rather
than relying on Schnorr batching alone.

\subsection{Tally Costs}
\label{sec:bench-tally}

The tally pipeline has two phases (Section~\ref{sec:phase-9-decryption}): 
serial-number decryption (used for deduplication) followed by aggregate 
decryption (used to recover the final tally). Both involve the same 
threshold-decryption primitive, applied to different ciphertexts.

\begin{table}[H]
\centering
\caption{Tally pipeline timings as a function of voter count $N$. 
``Serial decrypt per tallier'' is the wall-clock time for one tallier 
to perform her threshold-decryption share for all $N$ ballot serial 
numbers. ``Full pipeline'' is the end-to-end tally including 
aggregate decryption and BSGS lookup.}
\label{tab:bench-tally}
\small
\begin{tabular}{ccc}
\toprule
$N$ & Serial decrypt per tallier & Full pipeline \\
\midrule
16       & $1.832$ ms $\pm 0.3\%$  & $14.57$ ms $\pm 0.3\%$ \\
64       & $7.373$ ms $\pm 0.5\%$  & $50.96$ ms $\pm 0.5\%$ \\
256      & $29.36$ ms $\pm 0.4\%$  & $200.3$ ms $\pm 0.2\%$ \\
1{,}024  & $118.2$ ms $\pm 0.4\%$  & $888.0$ ms $\pm 0.5\%$ \\
4{,}096  & $474.4$ ms $\pm 0.4\%$  & $4.334$ s $\pm 0.3\%$ \\
16{,}384 & $2.004$ s $\pm 0.3\%$   & $31.78$ s $\pm 0.2\%$ \\
\bottomrule
\end{tabular}
\end{table}

\begin{figure}[H]
\centering
\includegraphics[width=0.85\textwidth]{g6_tally_breakdown.pdf}
\caption{Tally breakdown across $N$: per-tallier serial-decrypt cost
(blue) and full-pipeline cost (red) on the AWS run, log y-axis.
The full pipeline stays roughly an order of magnitude above the
serial step throughout, indicating aggregate decryption + BSGS lookup
dominate over the per-ballot partial decryption phase.
\figdata{thesis/results/timings.csv (group ends with \_tally)}}
\label{fig:g6-tally-breakdown}
\end{figure}

\begin{figure}[H]
\centering
\includegraphics[width=0.85\textwidth]{N16384_tally_full_pipeline_pdf.pdf}
\caption{Criterion sample-density plot for the full tally pipeline at
$N = 16{,}384$. The narrow peak around $31.8$ s confirms the reported
$\pm 0.2\%$ confidence in Table~\ref{tab:bench-tally}.
\figdata{thesis/results/criterion/N16384\_tally/full\_pipeline/report/pdf.svg}}
\label{fig:crit-tally-pipeline}
\end{figure}

Serial decryption scales linearly with $N$: each ballot has its own
encrypted serial number, and the tallier must partial-decrypt all of
them. From $N = 16$ to $N = 16{,}384$ (a $1024\times$ increase),
per-tallier serial decryption grows from $1.83$ ms to $2.00$ s
($\approx 1094\times$), close to the linear expectation; the slight
super-linearity at the top of the range is explained by cache pressure
from the larger ciphertext list on the $2$ MiB-per-vCPU AWS instance.

The full pipeline is roughly $8\times$ slower than per-tallier serial
decryption alone across the measured range, reflecting the additional
work of aggregate decryption (per-candidate ciphertexts), Lagrange
combination of tallier shares, and the BSGS lookup to convert
$g^{T_j}$ back into the integer $T_j$. At $N = 256$ the full pipeline
takes $200$ ms; at $N = 4{,}096$ it is $4.3$ s; at $N = 16{,}384$ it
is $31.8$ s. Linear extrapolation from the measured $N = 16{,}384$
point puts a single-instance tally at $N = 10^6$ near $\approx 32$
minutes; we mark this an extrapolation rather than a measurement
because the $N = 65{,}536$ attempt ran out of memory on this hardware
and could not be timed.

\section{End-to-End Election Simulation}
\label{sec:bench-lifecycle}

This section combines the per-phase costs into end-to-end timing for a
complete election across the measured voter scales. The configuration
is $\eta = 3$ authorities with threshold $t = 2$, default election
parameters, and $N \in \{16, 64, 256, 1{,}024, 4{,}096, 16{,}384\}$.

\begin{table}[H]
\centering
\caption{End-to-end election lifecycle timings (AWS run). ``Voter reg.''
is the per-voter registration cost (so total registration is roughly
$N \times$ that value). ``Ballot create'' and ``Ballot verify'' are
single-ballot times; total election ballot work is roughly
$N \times$ those values.}
\label{tab:bench-lifecycle}
\small
\begin{tabular}{cccccc}
\toprule
$N$ & DKG & Voter reg. & Ballot create & Ballot verify & Tally \\
\midrule
16       & $1.130$ ms & $119.2$ µs & $3.011$ ms & $1.622$ ms & $14.57$ ms \\
64       & $1.065$ ms & $117.0$ µs & $5.095$ ms & $2.160$ ms & $50.96$ ms \\
256      & $1.060$ ms & $116.2$ µs & $13.20$ ms & $4.266$ ms & $200.3$ ms \\
1{,}024  & $1.062$ ms & $117.6$ µs & $42.43$ ms & $10.64$ ms & $888.0$ ms \\
4{,}096  & $1.060$ ms & $118.0$ µs & $156.7$ ms & $37.35$ ms & $4.334$ s \\
16{,}384 & $1.091$ ms & $118.4$ µs & $668.2$ ms & $145.2$ ms & $31.78$ s \\
\bottomrule
\end{tabular}
\end{table}

Reading across rows, several patterns emerge:

\begin{itemize}
\item \textbf{DKG is constant in $N$}, as expected (it depends only on
$\eta$ and $t$): the measured value sits between $1.06$ ms and $1.13$ ms
across all six scales.
\item \textbf{Voter registration is constant in $N$}, because each voter's
work (one Pedersen commitment, one representation proof) is independent
of how many other voters there are; per-voter cost stays in a tight
$116$--$119$ µs band across the full range.
\item \textbf{Ballot creation grows with $N$}, dominated by the
commitment-set proof. From $N = 16$ to $N = 16{,}384$ (a $1024\times$
increase), ballot creation grows from $3.01$ ms to $668$ ms
($\approx 222\times$) --- much slower than $N$ itself but faster than
$\log N$ alone, consistent with the $O(N / \log N)$ amortized cost of
Pippenger multi-scalar multiplication in the prover.
\item \textbf{Ballot verification grows even more slowly}: from
$1.62$ ms at $N = 16$ to $145$ ms at $N = 16{,}384$ ($\approx 90\times$
over the same $1024\times$ scale). The verifier's MSM amortizes more
aggressively, and the bit-vector and Schnorr verifications are
constant in $N$.
\item \textbf{Tally grows linearly} (Section~\ref{sec:bench-tally}).
\end{itemize}

The aggregate picture: a boardroom-scale election with $N = 256$ takes
$\approx 1.1$ ms for setup, $30$ ms total voter registration
($256 \times 116$ µs), $3.4$ s total ballot construction
($256 \times 13.20$ ms; parallelizable across voters), $1.1$ s total
ballot verification ($256 \times 4.27$ ms; much less with batch
verification, see Table~\ref{tab:bench-batch}), and $0.20$ s for
tally. Total serialized wall-clock time is under $5$ seconds; in
practice ballots are constructed in parallel by different voters and
the protocol is dominated by verifier work, which on a single AWS vCPU
sits near a second.

Extending the same arithmetic to $N = 4{,}096$ gives $\approx 1.1$ ms
DKG, $0.48$ s total registration, $642$ s ($\approx 10.7$ min) total
serialized ballot construction, $153$ s ($\approx 2.6$ min) total
serialized ballot verification, and $4.3$ s tally; batch verification
compresses the verifier total significantly
(Section~\ref{sec:bench-verify}). At $N = 16{,}384$ those become
$3.04$ h serialized ballot construction ($16{,}384 \times 668.2$ ms),
$39.6$ min ballot verification ($16{,}384 \times 145.2$ ms), and
$31.8$ s tally --- numbers that argue for a parallel ballot pool and
for batched verification in any deployment beyond the boardroom
scale. Note that at this $N$ the Schnorr-style batch verification of
Section~\ref{sec:bench-verify} no longer compresses the verifier
total appreciably (Table~\ref{tab:bench-election-batch}), because the
commitment-set verifier inside each ballot is $O(N)$ and cannot be
folded into the batched multi-scalar multiplication.

\begin{figure}[H]
\centering
\includegraphics[width=0.85\textwidth]{g4_ballot_vs_n.pdf}
\caption{Per-ballot create and verify times as a function of $N$ on a
log-log scale (AWS run). Both grow super-linearly in raw $N$ but
sub-linearly in $\log_2 N$, consistent with the
$O(N / \log N)$ Pippenger amortization argued in the
text.
\figdata{thesis/results/timings.csv (group=N*\_ballot)}}
\label{fig:g4-ballot-vs-n}
\end{figure}

\begin{figure}[H]
\centering
\includegraphics[width=0.85\textwidth]{g4_full_pipeline_vs_n.pdf}
\caption{Tally full-pipeline cost as a function of $N$ on log-log
axes (AWS run). The straight slope confirms the linear-in-$N$
prediction across more than three decades of voter set size.
\figdata{thesis/results/timings.csv (group=N*\_tally, function=full\_pipeline)}}
\label{fig:g4-full-pipeline-vs-n}
\end{figure}

\begin{figure}[H]
\centering
\includegraphics[width=0.85\textwidth]{N16384_ballot_violin.pdf}
\caption{Criterion violin plot for the $N = 16{,}384$ ballot group on
the AWS run, showing the sample distributions for \texttt{create} and
\texttt{verify}. The fully-resolved distributions visualise the
$\pm 0.3\%$ / $\pm 0.2\%$ confidence figures reported in
Table~\ref{tab:bench-lifecycle}.
\figdata{thesis/results/criterion/N16384\_ballot/report/violin.svg}}
\label{fig:crit-N16384_ballot}
\end{figure}

\section{Proof Size Analysis}
\label{sec:bench-sizes}

Proof and ballot sizes are independent of runtime cost and were measured 
across the entire supported range up to $N = 1{,}048{,}576$ voters.

\subsection{Individual Proof Sizes}

The four Schnorr-style proofs have constant size, independent of any 
election parameter:

\begin{table}[H]
\centering
\caption{Sizes of the four Schnorr-style proofs.}
\label{tab:bench-schnorr-sizes}
\small
\begin{tabular}{lc}
\toprule
Proof system & Size \\
\midrule
Representation        & $96$ B \\
DLEQ                  & $96$ B \\
Encryption validity   & $128$ B \\
Serial validity       & $192$ B \\
\bottomrule
\end{tabular}
\end{table}

The bit-vector proof grows linearly with $k'$ (since each padded slot 
contributes its own committed bit):

\begin{table}[H]
\centering
\caption{Bit-vector proof size as a function of padded slot count $k'$.}
\label{tab:bench-bitvec-sizes}
\small
\begin{tabular}{cc|cc}
\toprule
$k'$ & Size & $k'$ & Size \\
\midrule
4    & $256$ B  & 16   & $640$ B  \\
5    & $288$ B  & 32   & $1{,}152$ B \\
8    & $384$ B  & 64   & $2{,}176$ B \\
\bottomrule
\end{tabular}
\end{table}

Even at $k' = 64$ slots, the bit-vector proof is just over $2$ KB --- 
far smaller than running $64$ separate bit proofs would be.

\subsection{Commitment-Set Proof Size}

This is the most interesting size measurement: the Bootle-Spark 
construction promises $O(\log N)$ proof size, and our measurements 
confirm this scaling clearly.

\begin{table}[H]
\centering
\caption{Commitment-set proof size as a function of voter count $N$, 
using binary ($n = 2$) digit decomposition. The size grows by exactly 
$320$ bytes for each doubling of $N$, confirming $O(\log N)$ scaling.}
\label{tab:bench-commset-sizes}
\small
\begin{tabular}{ccc|ccc}
\toprule
$N$ & $(n, m)$ & Size & $N$ & $(n, m)$ & Size \\
\midrule
$16$       & $(2, 4)$  & $672$ B  & $16{,}384$    & $(2, 14)$ & $2{,}272$ B \\
$64$       & $(2, 6)$  & $992$ B  & $65{,}536$    & $(2, 16)$ & $2{,}592$ B \\
$256$      & $(2, 8)$  & $1{,}312$ B & $262{,}144$ & $(2, 18)$ & $2{,}912$ B \\
$1{,}024$  & $(2, 10)$ & $1{,}632$ B & $1{,}048{,}576$ & $(2, 20)$ & $3{,}232$ B \\
$4{,}096$  & $(2, 12)$ & $1{,}952$ B & & & \\
\bottomrule
\end{tabular}
\end{table}

The pattern is clear and exact: each doubling of $N$ (one additional 
binary digit) adds exactly $320$ bytes to the proof. Going from $N = 16$ 
to $N = 10^6$ --- a $65{,}000\times$ increase in voter count --- proof 
size grows only $4.81\times$, from $672$ B to $3.23$ KB. This is the 
core efficiency property of the GK/Bootle construction in practice.

\subsection{Full Ballot Sizes}

A full Aura ballot includes the encryption validity proofs (one per 
slot), the bit-vector proof, the encrypted serial number with its 
validity proof, and the commitment-set proof. The total scales 
predominantly with $N$ (via the commitment-set proof), with smaller 
contributions from $k'$.

\begin{table}[H]
\centering
\caption{Total ballot size as a function of voter count $N$, with default 
election parameters ($k = 4$, $\ell_{\min} = 1$, $\ell_{\max} = 2$, 
giving $k' = 5$). Each doubling of $N$ adds $320$ bytes, the same as 
the commitment-set proof, confirming that proof is the dominant 
$N$-dependent term.}
\label{tab:bench-ballot-sizes}
\small
\begin{tabular}{cc|cc}
\toprule
$N$ & Ballot size & $N$ & Ballot size \\
\midrule
$16$       & $2{,}272$ B  & $16{,}384$    & $3{,}872$ B \\
$64$       & $2{,}592$ B  & $65{,}536$    & $4{,}192$ B \\
$256$      & $2{,}912$ B  & $262{,}144$   & $4{,}512$ B \\
$1{,}024$  & $3{,}232$ B  & $1{,}048{,}576$ & $4{,}832$ B \\
$4{,}096$  & $3{,}552$ B  & & \\
\bottomrule
\end{tabular}
\end{table}

A million-voter election produces ballots under $5$ KB each, despite
the $N = 10^6$ commitment list. For comparison, a typical TLS handshake
is several KB; an Aura ballot is the same order of magnitude as a
single web-page request.

\begin{figure}[H]
\centering
\includegraphics[width=0.9\textwidth]{g2_proof_sizes.pdf}
\caption{Proof and ballot byte sizes from \texttt{proof\_sizes.rs}
(log x-axis). The four Schnorr-style proofs and bit-vector proofs sit
at the low end ($96$--$2{,}176$ B); commitment-set proofs grow with
$\log N$; full ballot sizes scale from $2.2$ KB at $N = 16$ to under
$5$ KB at $N = 10^6$.
\figdata{thesis/results/sizes.csv (from raw/proof\_sizes.log)}}
\label{fig:g2-proof-sizes}
\end{figure}

\section{Comparison with the Original Paper's Estimates}
\label{sec:bench-vs-paper}

The original Aura paper (Section 4.1) gives only asymptotic complexity, 
not concrete numbers. We can therefore not directly compare wall-clock 
timings, but we can verify that the asymptotic claims hold in practice.

\paragraph{Claim 1: Commitment-set proof size is $O(\log N)$.}
The paper writes: \emph{``The size of this proof scales as 
$O(\log(N_{\text{voters}}))$ using the instantiation referenced.''} 
Our measurements (Table~\ref{tab:bench-commset-sizes}) confirm this 
exactly: every doubling of $N$ adds a constant $320$ bytes to the proof, 
which is $\Theta(\log N)$ growth. Going from $N = 16$ to $N = 10^6$, the 
proof grows $4.81\times$ while $N$ grows $65{,}000\times$ --- a 
near-perfect logarithmic relationship.

\paragraph{Claim 2: Verification scales as $O(N / \log N)$ via Pippenger.}
The paper writes: \emph{``While verification at first appears to scale 
as $O(N_{\text{voters}})$, the use of efficient multiscalar multiplication 
algorithms can reduce this complexity to 
$O(N_{\text{voters}} / \log(N_{\text{voters}}))$ with good constants.''} 
Our measurements show commitment-set verification growing from $230$
µs at $N = 16$ to $35.60$ ms at $N = 4{,}096$ (a $\approx 154\times$
increase for a $256\times$ increase in $N$). Naive $O(N)$ scaling
would predict a $256\times$ slowdown; the asymptotic
$O(N / \log_2 N)$ amortisation that Pippenger MSM theoretically
achieves predicts a $(4096/\log_2 4096) / (16/\log_2 16)
= (4096/12) / (16/4) \approx 85\times$ slowdown. The measured
$\approx 154\times$ sits cleanly between these two bounds, much closer
to the asymptotic ideal than to naive linear scaling --- the gap from
$85\times$ to $154\times$ reflects the constant factors and the
non-MSM work (transcript management, commitment computations,
serialization) that Pippenger amortisation cannot fold away.

\paragraph{Claim 3: Batch verification has constant marginal cost.}
The paper writes: \emph{``Verifier weighting of common group elements in 
the required multiscalar multiplication evaluation makes the marginal 
verification complexity constant, amortizing the overall cost across 
the batch.''} Our measurements (Table~\ref{tab:bench-batch}) confirm 
this clearly: at batch size $100$, marginal per-proof verification cost 
drops by $2$--$3\times$ across all three batchable proof systems. For 
encryption-validity specifically, batched cost is $60.29$ µs/proof at
batch $100$ (Table~\ref{tab:bench-batch}), versus $179.2$ µs/proof
individually --- exactly the amortization the paper predicts.

\paragraph{An empirical contribution: the $n$-ary trade-off.}
The Aura paper does not benchmark different choices of $n$ in the 
$n^m = N$ parameterization. Our measurements (Table~\ref{tab:bench-commset}) 
provide concrete data: the $n = 4$ variant approximately halves prover 
time while leaving verifier time unchanged. This is a meaningful 
practical observation, since prover time is the per-voter cost (paid 
once per ballot, on possibly resource-constrained client devices) while 
verifier time is the per-ballot infrastructure cost (paid centrally, 
amortizable). Choosing $n = 4$ over $n = 2$ is a Pareto improvement 
for client-side performance at no infrastructure cost.

\paragraph{Summary.}
All three of the paper's quantitative claims hold in our implementation. 
The asymptotic predictions are not just theoretical guarantees but 
practical reality: the implementation behaves exactly as the paper 
describes, with the additional empirical finding that the $n = 4$ digit 
base choice is strictly better than $n = 2$ on the tested platform.
% =============================================================
% CHAPTER 7 — DISCUSSION
% =============================================================
\chapter{Discussion}
\label{ch:discussion}

This chapter steps back from the implementation details of 
Chapter~\ref{ch:implementation} and the measurements of 
Chapter~\ref{ch:results} to consider what the work means in practice. 
We first ask whether Aura, as implemented, is viable for the kinds of 
elections the introduction motivated. We then catalog what Aura 
deliberately does not solve, what limitations remain in the present 
implementation, and which directions seem most promising for follow-up 
work.

\section{Practical Viability for Direct Democracy}
\label{sec:viability}

The benchmarks of Chapter~\ref{ch:results} let us answer the practical 
question concretely: at what election scale does Aura become workable, 
and what does each role pay?

\paragraph{Voter-side cost.}
The voter's wall-clock cost is dominated by the commitment-set
(eligibility) proof, which scales with the voter list size $N$.
Measured on the AWS instance (2 vCPU Xeon Platinum 8488C): at
$N = 256$ a ballot takes $13.2$ ms to construct, at $N = 1{,}024$
$42.4$ ms, at $N = 4{,}096$ $157$ ms, and at $N = 16{,}384$
$668$ ms. Extrapolating linearly in $N$ from the measured
$N = 16{,}384$ point places $N = 10^6$ ballot construction near
$\approx 41$ s on the same single-vCPU hardware; we mark this an
extrapolation rather than a measurement because the $N = 65{,}536$
attempt exhausted the $3.7$ GiB instance memory and could not be
timed. For a phone-grade CPU we would expect proof times perhaps
$5$--$10\times$ longer (an estimate rather than a measurement; no
ARM benchmarks were run for this thesis). This places Aura
comfortably in the practical range for elections up to mid-six-figure
voter counts (university elections, professional-society votes,
municipal referenda) on commodity client devices, and within reach
for city- to country-scale elections if voters are willing to wait
a minute or two for a ballot to be prepared --- comparable to the
time spent in a physical polling booth.

The empirical observation that $n = 4$ approximately halves prover 
time at no verifier cost (Section~\ref{sec:bench-commset}) is 
particularly valuable here, since prover work is precisely the 
client-side cost paid on the most resource-constrained device in the 
system. A deployment that chooses $n = 4$ as the default --- as the 
present implementation can do whenever $N$ is a power of $4$ --- 
gets the improvement essentially for free.

\paragraph{Verifier-side cost.}
Verification is the cost paid by the public bulletin board, by
auditors, and by any party who wishes to confirm the election outcome.
Per-ballot verification at $N = 256$ takes $4.27$ ms; at $N = 4{,}096$
$37.4$ ms; at $N = 16{,}384$ $145$ ms (all measured on the AWS
instance). Crucially, the per-ballot cost is amortizable across
batches: encryption-validity verification gains close to $3\times$ at
batch size $100$ (Section~\ref{sec:bench-verify}), and the same
applies to the representation and DLEQ proofs used elsewhere in the
protocol; for the encryption-validity proof inside a ballot, batching
delivers close to $3\times$ amortisation
(Section~\ref{sec:bench-verify}). For a full election of
$10{,}000$ ballots at $N = 16{,}384$, however, the measured per-ballot
verify of $145$ ms does \emph{not} compress to $\approx 50$ ms with
$k = 100$ batching --- the commitment-set verifier inside each ballot
is $O(N)$ and cannot be folded into the batched multi-scalar
multiplication, so the verifier total stays near
$10{,}000 \times 145$ ms $\approx 24$ minutes single-threaded. A
practical deployment closes this gap with the parallel verifier pool
that any real bulletin board would deploy; on the AWS instance used
here ($2$ vCPU) the speedup ceiling is modest, but on
$16$--$32$-vCPU verifier hardware the total drops below the
$2$--$3$ minute audit-window typically targeted for elections of this
size.

\paragraph{Authority-side cost.}
The talliers do constant-time work in the DKG phase
($\approx 1.1$ ms for $\eta = 3$, $t = 2$, independent of $N$) and
linear work in the tally phase. Tally-pipeline cost grows roughly
linearly with $N$: $14.6$ ms at $N = 16$, $200$ ms at $N = 256$,
$4.33$ s at $N = 4{,}096$, $31.8$ s at $N = 16{,}384$. Linear
extrapolation from the measured $N = 16{,}384$ row places a single-
instance tally at $N = 10^6$ near $\approx 32$ minutes --- still the
most $N$-dependent phase of the protocol, but it runs once at the end
of the election on dedicated infrastructure, and that figure should be
read as an extrapolation rather than a measurement: the $N = 65{,}536$
attempt did not complete on the $3.7$ GiB AWS instance.

\paragraph{Verdict.} 
For boardroom and mid-scale elections (up to roughly $10^4$ voters), 
Aura on the present implementation is straightforwardly practical: 
voters wait under a second, verifiers complete in seconds, talliers 
finish in milliseconds. For larger scales (up to $10^6$ voters), the 
protocol remains tractable but voter-side wait times begin to matter 
on weak hardware --- this is the regime where the future-work 
directions discussed in Section~\ref{sec:future-work} become most 
relevant.

\section{What Aura Does Not Solve}
\label{sec:limits}

It is important to be honest about what Aura, even with this 
implementation, does \emph{not} provide. Several long-standing 
challenges in electronic voting are deliberately out of scope.

\paragraph{Receipt-freeness.} 
Aura does not prevent a willing voter from proving how she voted. A 
voter who chooses to retain her ballot's encryption randomness 
$\{\rho_{i, j}\}$ can later show a third party that she encrypted 
specific values into specific slots, which constitutes a verifiable 
receipt. This makes Aura unsuitable for jurisdictions where 
receipt-freeness is mandated by law, and limits its protection 
against vote-buying schemes that rely on voters voluntarily 
proving their ballot contents. True receipt-freeness 
generally requires a stronger trust model (a designated verifier, 
re-encryption shuffle, or hardware-protected randomness) than Aura 
adopts.

\paragraph{Vote-buying mitigation beyond coercion resistance.} 
Aura's coercion resistance covers the specific case of a voter under 
duress who later regains autonomy and re-votes. It does not cover the 
weaker but more pervasive case of voluntary vote-selling, where a 
voter willingly reveals her receipt for compensation. The 
JCJ-credentials approach~\cite{jcj} mitigates this with anonymous 
credentials that voters can fake under duress; integrating such a 
mechanism into Aura is non-trivial and is a natural follow-up question.

\paragraph{Sybil resistance and voter authentication.} 
Aura assumes a list $L_{\text{voters}}$ of authorized signing keys 
distributed by an organizer. The mechanism that establishes this list 
--- proving that each registered voter is a real, unique, eligible 
person --- is outside the protocol. A real deployment would have to 
rely on existing identity infrastructure (national ID systems, 
attestations from independent registrars, or web-of-trust schemes) 
to populate $L_{\text{voters}}$. The cryptographic guarantees Aura 
provides are conditional on this list being correct.

\paragraph{Bulletin-board availability.} 
Aura assumes a publicly readable, append-only bulletin board where 
ballots are posted and on which observers can run \texttt{VerifyBallot}. 
The protocol does not specify how the bulletin board is implemented, 
nor how to defend against an adversarial bulletin board that silently 
drops ballots. A blockchain bulletin board would address availability 
at the cost of latency and storage; a centralized bulletin board with 
replicated audit logs would be cheaper but trust-bound.

\paragraph{Network anonymity.} 
Aura provides cryptographic ballot anonymity but not network-level 
anonymity. A voter who connects directly from a known IP address to 
post her ballot reveals when she voted, even if not what she voted for. 
A real deployment would route ballot submissions through a 
mixnet or Tor.

\section{Implementation Limitations}
\label{sec:impl-limits}

The implementation described in Chapter~\ref{ch:implementation} is a 
faithful reference for the protocol but is not yet production-grade. 
We catalog the limitations explicitly so that any party wishing to 
deploy or build on this code can address them deliberately.

\paragraph{Variable-time arithmetic.} 
The implementation uses 
\texttt{curve25519-dalek}'s \texttt{vartime\_multiscalar\_mul} for all 
prover and verifier multi-scalar multiplications. This is the standard 
choice for verifier code (where inputs are public) but is incorrect 
for prover code in adversarial settings, where timing channels could 
leak the secret index $\ell$ or seed $s_i$. Production deployment 
would require switching prover-side MSM to constant-time variants 
(at a roughly $2\times$ speed cost) and auditing every other 
arithmetic operation for timing-channel safety. The 
\texttt{curve25519-dalek} crate provides constant-time scalar arithmetic, 
so this is largely a matter of replacing the variable-time MSM call 
with the constant-time one and re-running the benchmarks.

\paragraph{No formal security audit.} 
The implementation has not been independently audited. Test coverage 
($57$ unit tests plus $5$ integration tests, all passing under 
\texttt{cargo clippy --deny warnings}) catches functional regressions 
but cannot substitute for a third-party cryptographic review. 
Production deployment would require an audit by an established 
firm with cryptographic-implementation experience. A formal-methods 
audit (machine-checked proof that the Rust code matches the 
$\Sigma$-protocol specifications of the paper) would be a stronger 
guarantee but a much heavier engineering investment.

\paragraph{No persistent bulletin board.} 
The current implementation simulates the bulletin board as an 
in-memory \texttt{Vec} of accepted ballots. A real deployment would 
need integration with a persistent, tamper-evident store (a 
blockchain, a transparency log, or a replicated append-only database) 
together with the gossip and ordering protocol that surrounds it. 
This is a substantial systems engineering effort, separate from the 
cryptographic work documented here.

\paragraph{No client SDK or user interface.} 
Voters, talliers, and organizers all interact with the protocol 
through Rust function calls. Real deployment would require a client 
application (mobile app, browser extension, or web frontend), a UX 
flow for ballot construction and re-voting, an authentication layer 
that ties signing keys to voter identities, and an operations 
dashboard for talliers. None of this is in scope for the present 
thesis.

\paragraph{Measurement ceiling at $N = 65{,}536$.}
The runtime measurements in Chapter~\ref{ch:results} cover
$N \in \{16, 64, 256, 1{,}024, 4{,}096, 16{,}384\}$ on the AWS
instance described in Section~\ref{sec:bench-methodology}. The
attempt at $N = 65{,}536$ exhausted the instance's $3.7$ GiB of RAM
during ballot construction and produced no timing data; we treat
that configuration as a soft ceiling for the hardware used here.
Larger-$N$ measurements are not blocked by the implementation
(proof-size measurements extend up to $N = 10^6$ on the same
binary), only by per-instance memory headroom. A larger AWS
instance type or a workstation with $\ge 16$ GiB of RAM should
reach $N = 10^6$ for the runtime benchmarks; that work is a
roughly half-day exercise but was not in scope for the present
thesis.

\paragraph{Threats to benchmark validity.}
A defensive reading of Chapter~\ref{ch:results} should keep in mind
what was \emph{not} measured. (i) All runtime benchmarks were taken
on a single AWS instance with $2$ vCPU; multi-vCPU scaling is
unmeasured, so the parallel-verifier claims in the viability
discussion (Section~\ref{sec:viability}) are projections, not
measurements. (ii) No mobile-/ARM-class CPU was benchmarked; the
``$5$--$10\times$ slower on a phone'' framing in
Section~\ref{sec:viability} is an estimate from public benchmarks
of \texttt{curve25519-dalek}, not a number we measured for this
thesis. (iii) The hand-rolled vs library-backed comparison
(Appendix~\ref{app:lib-comparison}) was run on a separate laptop
because the \texttt{lib-proofs} Cargo feature was not exercised on
the AWS instance; absolute timings between the body and the
appendix are not directly comparable, only the
hand-rolled-vs-library ratios are. (iv) No heap/RSS memory
profiling was performed, so the $N = 65{,}536$ ceiling is observed
empirically rather than attributed to a specific allocation. (v)
The implementation uses variable-time MSM throughout; constant-time
prover work (required for adversarial-timing settings) would slow
prover measurements by roughly $2\times$ and the benchmarks here
do not reflect that overhead.

\section{Future Work}
\label{sec:future-work}

We organize future directions roughly in order of how directly they 
build on the present thesis.

\paragraph{Standard-tier and full-tier benchmarks.} 
Most immediately, running the existing benchmark infrastructure at 
\emph{standard} and \emph{full} tiers would extend wall-clock 
measurements up to $N = 10^6$, validating the extrapolations in 
Section~\ref{sec:viability} with direct data. This is a 
one-engineering-day task that would substantially strengthen the 
practical-viability claims of the thesis.

\paragraph{Hand-rolled vs library-backed comparison.} 
Producing the empirical comparison between Aura's hand-rolled 
$\Sigma$-protocols and equivalent off-the-shelf Rust ZK libraries 
(\texttt{zkp}, \texttt{sigma-protocols}, \texttt{bulletproofs}, 
\texttt{one-of-many-proofs}) would provide useful guidance for 
implementers choosing between paper-faithful custom code and 
library-friendly approximations. The infrastructure for this 
comparison is in place behind the \texttt{lib-proofs} feature flag; 
running the comparison and writing up the statement-fidelity 
trade-offs would yield a publishable empirical contribution.

\paragraph{Constant-time prover and security audit.} 
Switching prover-side MSM from variable-time to constant-time is a 
mechanical change that would close the most pressing implementation 
gap before any real deployment. A subsequent third-party security 
audit, ideally combined with formal verification of the Rust code 
against the protocol's $\Sigma$-protocol specifications, would 
establish a level of assurance comparable to other deployed 
e-voting systems.

\paragraph{SNARK-based eligibility proofs with Merkle anchoring.} 
A natural research extension is to replace Aura's flat-list 
Groth-Kohlweiss commitment-set proof with a Merkle-tree-anchored 
zk-SNARK eligibility proof, in the style of Zcash 
Sapling~\cite{zcash-sapling}, Tornado Cash, or 
Semaphore~\cite{semaphore}. The voter list would compress from $O(N)$ 
to a single $O(1)$ Merkle root on chain, eligibility proofs would 
shrink from $O(\log N)$ to constant size (~192 bytes for Groth16), 
and verification would become constant-time per ballot --- significant 
improvements at large scale ($N > 10^6$). The trade-off is the choice 
of proof system. Groth16 offers the smallest proofs and fastest 
verification but requires a per-circuit trusted setup ceremony, 
directly conflicting with Aura's stated trust-free setup design 
goal. Transparent alternatives such as Halo2 with FRI commitments, 
STARKs, or folding schemes like Nova~\cite{nova} preserve setup-freeness 
at the cost of larger proofs and heavier prover work. A direct 
empirical comparison of these variants against the current Aura would 
quantify the practical trade-offs and clarify which design is best 
suited to which deployment scale. Adjacent ideas in this direction 
include incremental Merkle trees~\cite{zcash-sapling} for 
late-registration support during voting, and lightweight client 
membership proofs for auditors who do not want to download the full 
voter list.

\paragraph{Integration with a real bulletin board.} 
Connecting the protocol to a persistent bulletin board --- a 
blockchain, a transparency log, or a Tor-anchored replicated store 
--- would move the implementation from research-grade to deployment-grade. 
The cryptographic interface is well-defined; the engineering effort 
is in the surrounding consensus, gossip, and storage layers.

\paragraph{Stronger coercion and receipt-freeness mechanisms.} 
Integrating JCJ-style anonymous credentials~\cite{jcj} would extend 
Aura's weak coercion-resistance to a stronger property, allowing 
voters under sustained duress (not just transient duress) to cast 
fake ballots that are silently discarded at tally time. This is a 
substantial protocol modification rather than an implementation 
extension and would warrant its own thesis-scale project.

\paragraph{Faster elliptic curves.} 
Ristretto255 was chosen for its constant-time properties and its 
mature ecosystem in Rust. Newer pairing-friendly or AVX-optimized 
curves (BLS12-381 with batch operations, or recent assembly-optimized 
implementations of Curve25519-derived constructions) could yield 
$2$--$5\times$ speedups in the most expensive operations. This is 
a direct engineering optimization; the protocol would not change.

\section{Concluding Remarks}
\label{sec:concluding-remarks}

This thesis has addressed the question of whether the Aura private 
voting protocol of Jivanyan and Feickert~\cite{aura2022} is practical 
in implementation, and what concrete trade-offs surface when it is 
built end-to-end on modern hardware. The answer is yes, with caveats 
worth stating plainly.

Aura's design priorities --- trust-free setup, simple proof systems, 
universal verifiability, and weak coercion resistance --- are real, 
well-motivated gaps in the existing electronic-voting landscape. 
Helios, ElectionGuard, and the various deployed systems discussed in 
Chapter~\ref{ch:literature} each compromise on at least one of these 
properties, and Aura's specific combination of guarantees is not 
redundant with prior work.

This thesis delivers what is, to our knowledge, the first fully 
open-source implementation of the protocol, faithful to the original 
paper at every transcript byte and generator choice. The benchmarks 
of Chapter~\ref{ch:results} show that Aura is straightforwardly 
practical for boardroom-scale elections (under two seconds of 
cumulative cryptographic work for $N \le 10^3$), workable for 
municipal-scale elections (under a minute for $N \le 10^5$), and 
within reach for national-scale elections with the engineering 
investments outlined above. The asymptotic claims of the original 
paper hold exactly in our measurements; as an additional empirical 
contribution, we identified that the $n = 4$ digit-base in the 
Bootle-Spark commitment-set proof is approximately twice as fast on 
the prover side as the binary $n = 2$ default, at no cost to the 
verifier --- a Pareto improvement that the paper did not benchmark.

The worked numerical example of Chapter~\ref{ch:byhand}, the 
protocol description of Chapter~\ref{ch:aura}, the cryptographic 
primitives of Chapter~\ref{ch:primitives}, and the implementation 
documentation of Chapter~\ref{ch:implementation} should together be 
sufficient for a careful reader to understand, audit, and build on 
Aura without prior cryptographic background beyond the basics. The 
code is licensed BSD-3-Clause and is available at 
\url{https://github.com/nerses-asaturyan/aura}.

The most important next steps, in priority order, are: (1) extending 
the wall-clock benchmarks from the \emph{quick} tier to the 
\emph{standard} and \emph{full} tiers, validating the extrapolations 
in this thesis with direct measurement at $N$ up to $10^6$; (2) 
replacing variable-time prover arithmetic with constant-time variants 
and securing a third-party cryptographic audit, closing the most 
pressing gap between research-grade and deployment-grade implementation; 
and (3) investigating Merkle-anchored zk-SNARK eligibility proofs as 
a route to constant-size ballot proofs at national scale, with 
transparent proof systems (Halo2 with FRI, STARKs, Nova) preserving 
Aura's trust-free setup property. Together these extensions would 
take Aura from a paper-faithful research artifact to a deployable 
e-voting infrastructure suitable for the direct-democracy applications 
that motivated this work.

% =============================================================
% APPENDIX B — Library-backed proof comparison (laptop run)
% =============================================================
\appendix
\chapter{Library-backed proof comparison}
\label{app:lib-comparison}

\section*{A.1 Setup}

The benchmarks reported in Chapter~\ref{ch:results} were measured on the
AWS instance described in Section~\ref{sec:bench-methodology}. The
hand-rolled vs library-backed comparison in this appendix was measured
on a separate, pre-existing local machine because the
\texttt{lib-proofs} Cargo feature pulls in additional dependencies
(\texttt{bulletproofs} git-main, \texttt{zkp}, \texttt{sigma-rs}) that
were not part of the AWS environment. The laptop configuration was:
12th-generation Intel Core i7-12700H (10 cores, 20 threads, $24$ MiB
L3), $16$ GiB RAM, Linux kernel $6.6.87$ inside WSL2 (Ubuntu $22.04$
host), \texttt{rustc 1.93.1} / \texttt{cargo 1.93.1}, release profile
with default LTO. The randomness seed (\texttt{ChaCha20Rng} state
$20260412$) is identical to the AWS run. The full environment dump
lives in \texttt{thesis/results/env\_laptop.txt}; the raw bench log is
\texttt{thesis/results/raw/proof\_primitives\_lib.log} and the
underlying Criterion tree is at
\texttt{thesis/results/criterion\_g5\_laptop/}.

\section*{A.2 Per-proof comparison}

Table~\ref{tab:g5_handrolled_vs_lib} pairs the hand-rolled
implementation against its library counterpart for each batchable proof
system, in both \texttt{prove} and \texttt{verify} modes. The
right-most column reports the ratio
$t_{\text{library}} / t_{\text{hand-rolled}}$: values close to $1$
indicate parity; $> 1$ means the library is slower.

\begin{table}[h]
  \centering
  \caption{Hand-rolled vs library-backed prove/verify (statement-fidelity caveats noted in Section X).}
  \label{tab:g5_handrolled_vs_lib}
  \begin{tabular}{lrrrrr}
    \toprule
    Group & Op & Scale & Hand-rolled & Library & lib / hand-rolled \\
    \midrule
bitvec\_compare & prove & 16 & 932.4 $\mu$s & 5.636 ms & 6.04× \\
bitvec\_compare & prove & 64 & 2.961 ms & 20.25 ms & 6.84× \\
bitvec\_compare & verify & 16 & 569.1 $\mu$s & 1.503 ms & 2.64× \\
bitvec\_compare & verify & 64 & 1.860 ms & 4.609 ms & 2.48× \\
dleq\_compare & prove & — & 126.8 $\mu$s & 119.1 $\mu$s & 0.94× \\
dleq\_compare & verify & — & 207.3 $\mu$s & 219.4 $\mu$s & 1.06× \\
encval\_compare & prove & — & 184.6 $\mu$s & 243.3 $\mu$s & 1.32× \\
encval\_compare & verify & — & 282.7 $\mu$s & 196.4 $\mu$s & 0.69× \\
rep\_compare & prove & 1 & 63.21 $\mu$s & 104.9 $\mu$s & 1.66× \\
rep\_compare & prove & 2 & 80.40 $\mu$s & 144.8 $\mu$s & 1.80× \\
rep\_compare & prove & 4 & 116.8 $\mu$s & 214.4 $\mu$s & 1.84× \\
rep\_compare & verify & 1 & 104.4 $\mu$s & 87.49 $\mu$s & 0.84× \\
rep\_compare & verify & 2 & 123.4 $\mu$s & 111.6 $\mu$s & 0.90× \\
rep\_compare & verify & 4 & 162.1 $\mu$s & 155.8 $\mu$s & 0.96× \\
serval\_compare & prove & — & 295.2 $\mu$s & 363.4 $\mu$s & 1.23× \\
serval\_compare & verify & — & 384.0 $\mu$s & 300.2 $\mu$s & 0.78× \\
set\_compare & prove & N=1024 & 64.87 ms & 596.1 ms & 9.19× \\
set\_compare & prove & N=256 & 18.97 ms & 123.5 ms & 6.51× \\
set\_compare & verify & N=1024 & 14.09 ms & 72.22 ms & 5.13× \\
set\_compare & verify & N=256 & 4.017 ms & 20.36 ms & 5.07× \\
    \bottomrule
  \end{tabular}
\end{table}


The results are mixed by design. The hand-rolled stack carries minimal
abstraction overhead and is consistently competitive with the library
implementations on the simple Schnorr-style relations
(\texttt{rep\_compare}, \texttt{dleq\_compare},
\texttt{encval\_compare}, \texttt{serval\_compare}). On the heavier
proof systems (\texttt{bitvec\_compare}, \texttt{set\_compare}), the
library variants use mature inner-product / one-of-many primitives
that exploit further optimisations the hand-rolled code does not
attempt, so they are typically faster at the cost of larger code
footprint and an additional dependency surface.

\section*{A.3 Figures}

\begin{figure}[H]
\centering
\includegraphics[width=0.9\textwidth]{g5_per_proof_compare_prove.pdf}
\caption{Hand-rolled vs library-backed prove times, per proof system,
on the laptop reference run. Log y-axis. Heights are absolute
microseconds; the relative gap is the value of interest.
\figdata{thesis/results/criterion\_g5\_laptop/*\_compare/prove/}}
\label{fig:g5-prove}
\end{figure}

\begin{figure}[H]
\centering
\includegraphics[width=0.9\textwidth]{g5_per_proof_compare_verify.pdf}
\caption{Hand-rolled vs library-backed verify times, per proof system,
on the laptop reference run. Log y-axis.
\figdata{thesis/results/criterion\_g5\_laptop/*\_compare/verify/}}
\label{fig:g5-verify}
\end{figure}

\section*{A.4 Threats to portability}

Absolute timings in this appendix are not directly comparable to the
AWS numbers in the body of the thesis: the laptop's i7-12700H runs at
substantially higher single-core frequency, has $\approx 4\times$ more
RAM, and is not virtualised. The appendix focuses on the
\emph{ratio} between hand-rolled and library variants, which is
expected to be hardware-portable to first order because both
implementations exercise the same Ristretto255 group operations in the
same proportions. A future run that exercises the
\texttt{lib-proofs} feature on the AWS instance would let us make
absolute comparisons single-hardware; that work is left for future
iterations of the codebase.

% =============================================================
% BIBLIOGRAPHY
% =============================================================
\addcontentsline{toc}{chapter}{Bibliography}
\begin{thebibliography}{99}

\bibitem{aura2022}
Jivanyan A., Feickert A.,
\textit{Aura: private voting with reduced trust on tallying authorities},
Cryptology ePrint Archive, Paper 2022/543, 2022.
\url{https://eprint.iacr.org/2022/543}

\bibitem{helios}
Adida B.,
\textit{Helios: Web-based Open-Audit Voting},
USENIX Security Symposium, 2008, pp.~335--348.

\bibitem{russian2021}
Maksimov V., Smyshlyaev S.,
\textit{Russian Federal Remote E-voting Scheme of 2021},
ePrint Archive, Paper 2021/1454, 2021.
\url{https://eprint.iacr.org/2021/1454}

\bibitem{cryptoreview}
Bernhard D., Cortier V., Galindo D., Pereira O., Warinschi B.,
\textit{A comprehensive analysis of game-based ballot privacy definitions},
IEEE Symposium on Security and Privacy, 2015, pp.~499--516.

\bibitem{blockchainreview}
Pawlak M., Guziur J., Poniszewska-Mara\'{n}da A.,
\textit{Towards the blockchain technology for system voting process},
Procedia Computer Science, vol.~159, 2019, pp.~2293--2301.

\bibitem{merkle}
Merkle R.~C.,
\textit{A digital signature based on a conventional encryption function},
CRYPTO, 1987, pp.~369--378.

\bibitem{groth-kohlweiss}
Groth J., Kohlweiss M.,
\textit{One-out-of-many proofs: Or how to leak a secret and spend a coin},
EUROCRYPT, 2015, pp.~253--280.

\bibitem{pedersen}
Pedersen T.~P.,
\textit{Non-interactive and information-theoretic secure verifiable secret sharing},
CRYPTO, 1991, pp.~129--140.

\bibitem{shamir}
Shamir A.,
\textit{How to share a secret},
Communications of the ACM, vol.~22, no.~11, 1979, pp.~612--613.

\bibitem{elgamal}
ElGamal T.,
\textit{A public key cryptosystem and a signature scheme based on discrete logarithms},
IEEE Transactions on Information Theory, vol.~31, no.~4, 1985, pp.~469--472.

\bibitem{schnorr}
Schnorr C.~P.,
\textit{Efficient signature generation by smart cards},
Journal of Cryptology, vol.~4, no.~3, 1991, pp.~161--174.

\bibitem{fiatshamir}
Fiat A., Shamir A.,
\textit{How to prove yourself: Practical solutions to identification and signature problems},
CRYPTO, 1986, pp.~186--194.

\bibitem{cramer-damgard}
Cramer R., Damg\aa rd I., Schoenmakers B.,
\textit{Proofs of partial knowledge and simplified design of witness hiding protocols},
CRYPTO, 1994, pp.~174--187.

\bibitem{feldman}
Feldman P.,
\textit{A practical scheme for non-interactive verifiable secret sharing},
FOCS, 1987, pp.~427--438.

\bibitem{bellare-rogaway}
Bellare M., Rogaway P.,
\textit{Random oracles are practical: A paradigm for designing efficient protocols},
ACM Conference on Computer and Communications Security, 1993, pp.~62--73.

\bibitem{bootle-bit-vector}
Bootle J., Cerulli A., Chaidos P., Ghadafi E., Groth J., Petit C.,
\textit{Short accountable ring signatures based on DDH},
ESORICS, 2015, pp.~243--265.

\bibitem{lelantus-spark}
Jivanyan A., Feickert A.,
\textit{Lelantus Spark: Secure and flexible private transactions},
Cryptology ePrint Archive, Paper 2021/1173, 2021.
\url{https://eprint.iacr.org/2021/1173}

\bibitem{jcj}
Juels A., Catalano D., Jakobsson M.,
\textit{Coercion-resistant electronic elections},
ACM Workshop on Privacy in the Electronic Society, 2005, pp.~61--70.

\bibitem{zcash-sapling}
Hopwood D., Bowe S., Hornby T., Wilcox N.,
\textit{Zcash protocol specification, Version 2024.7.0 [Sapling/Orchard]},
Electric Coin Company, 2024.
\url{https://zips.z.cash/protocol/protocol.pdf}

\bibitem{semaphore}
Whitehat B., Mueller J., Beregszaszi A.,
\textit{Semaphore: Anonymous signaling on Ethereum},
Ethereum Foundation, 2019.
\url{https://semaphore.pse.dev}

\bibitem{nova}
Kothapalli A., Setty S., Tzialla I.,
\textit{Nova: Recursive zero-knowledge arguments from folding schemes},
CRYPTO, 2022, pp.~359--388.

\end{thebibliography}

% =============================================================
% APPENDICES (optional)
% =============================================================
\end{document}